% FST-NS_LogDistanceIntegrability_v1_de.tex
\documentclass[12pt,a4paper]{article}

\usepackage[utf8]{inputenc}
\usepackage[T1]{fontenc}
\usepackage{lmodern}
\usepackage{microtype}
\usepackage{amsmath,amssymb,amsthm}
\usepackage{enumitem}
\usepackage{xcolor}
\usepackage{geometry}
\usepackage{graphicx}
\geometry{a4paper, margin=2.5cm}
\usepackage{hyperref}

\hypersetup{
    unicode=true,
    colorlinks=true,
    linkcolor=blue!70!black,
    citecolor=green!50!black,
    urlcolor=blue!70!black,
    pdftitle={Log-Distanz-Integrierbarkeit für dissipative Strömungen},
    pdfauthor={Lukas Geiger},
    pdfsubject={BV-Selektionen auf fraktalen Attraktoren},
    pdfkeywords={Log-Distanz-Integrierbarkeit, trajektorienweise log-Lipschitz, beschränkte Variation, dissipative PDEs, Navier-Stokes}
}
\renewcommand{\theHequation}{de.\thesection.\arabic{equation}}
\providecommand*{\HyperDestNameFilter}[1]{de.#1}
\setlength{\emergencystretch}{4em}

\theoremstyle{definition}
\newtheorem{definition}{Definition}[section]
\theoremstyle{plain}
\newtheorem{theorem}[definition]{Satz}
\newtheorem{lemma}[definition]{Lemma}
\newtheorem{proposition}[definition]{Proposition}
\newtheorem{corollary}[definition]{Korollar}
\newtheorem{conjecture}[definition]{Vermutung}
\theoremstyle{remark}
\newtheorem{remark}[definition]{Bemerkung}
\newtheorem{example}[definition]{Beispiel}
\renewcommand{\theHdefinition}{de.\thesection.\arabic{definition}}

\renewcommand{\proofname}{Beweis}

\newcommand{\Var}{\mathrm{Var}}
\newcommand{\dist}{\mathrm{dist}}
\newcommand{\AC}{\mathrm{AC}}
\newcommand{\BV}{\mathrm{BV}}
\newcommand{\id}{\mathrm{id}}

\title{\textbf{Log-Distanz-Integrierbarkeit für dissipative Strömungen:\\
BV-Selektionen auf fraktalen Attraktoren}}
\author{Lukas Geiger\\
\small Unabhängiger Forscher, Bernau, Deutschland\\
\small ORCID: \href{https://orcid.org/0009-0005-7296-1534}{0009-0005-7296-1534}%
\thanks{KI-Werkzeuge (Claude von Anthropic) wurden für mathematische Formalisierung, Literaturverifikation und redaktionelle Verfeinerung eingesetzt. Alle mathematischen Inhalte wurden vom Autor entwickelt, verifiziert und freigegeben.}}
\date{Mai 2026}

\begin{document}
\pagenumbering{gobble}
\maketitle

\begin{center}
\small
\textbf{MSC 2020:} 37L30, 47H09, 46B20, 35Q30, 26A45\\
\textbf{Schlüsselwörter:} Nächstpunkt-Projektion, log-Lipschitz-Regularität,
beschränkte Variation, Attraktorgeometrie, dissipative PDEs, Navier--Stokes
\end{center}

\begin{abstract}
Wir führen einen Regularitätsrahmen für Nächstpunkt-Projektionen
auf kompakte Attraktoren dissipativer Semiflows in Hilberträumen ein.
Das klassische Lipschitz-Projektionstheorem (Federer, 1959) erfordert
positive Reichweite -- eine Bedingung, die mit fraktaler
Attraktorgeometrie unvereinbar ist.  Wir ersetzen sie durch eine
\emph{trajektorienweise log-Lipschitz}-Bedingung (TLL), die den
Projektionsmodul auf der geometrisch natürlichen Skala des
Attraktorabstands statt der Partitionsfeinheit kontrolliert.
In Kombination mit einer \emph{Log-Distanz-Integrierbarkeits}-Bedingung
(LDI) ergibt dies Selektionen beschränkter Variation entlang absolut
stetiger Trajektorien -- genau die Regularität, die für
Gronwall-artige Argumente in bedingten Regularitätskriterien
benötigt wird.  Der Rahmen gilt direkt für dissipative Semiflows
mit kompaktem Attraktor endlicher fraktaler Dimension; für die
3D-Navier-Stokes-Gleichungen ist er als bedingte Schnittstelle für
einen starken oder vorab gewählten Trajektorienattraktor-Rahmen zu
lesen. Reaktions-Diffusions-Systeme und gedämpfte Wellengleichungen
liefern sauberere Testklassen.  Zwei verifizierbare hinreichende
Bedingungen für LDI werden identifiziert: beschränkte Variation
der Log-Distanz-Funktion und potenzgesetzliche Sublevel-Zeitkontrolle.
\end{abstract}

\clearpage
\pagenumbering{arabic}

% ==================================================================
\section{Einleitung}
% ==================================================================

Sei $H$ ein separabler Hilbertraum und $\mathcal{A} \subset H$ eine
kompakte Menge -- der globale Attraktor eines dissipativen Semiflows
$\{S(t)\}_{t \ge 0}$.  Eine grundlegende Frage in der qualitativen
Theorie dissipativer partieller Differentialgleichungen ist die
Regularität der Nächstpunkt-Projektion
\[
  P(x) \;\in\; \arg\min_{a \in \mathcal{A}} \|x - a\|_H
\]
als Funktion der Zeit entlang einer Trajektorie $u(t) = S(t)\,u_0$.
Ist die Abbildung $t \mapsto P(u(t))$ absolut stetig (AC), so
schließen sich Gronwall-artige Energieabschätzungen automatisch;
ist sie lediglich messbar, bleiben kritische Fehlerterme
unkontrolliert.

Die klassische Theorie (Federer~\cite{Federer1959}) garantiert
Lipschitz-Stetigkeit der Nächstpunkt-Projektion auf der tubularen
Umgebung von Mengen mit \emph{positiver Reichweite}.  Attraktoren
dissipativer PDEs haben jedoch generisch nicht-ganzzahlige fraktale
Dimension und Reichweite null
(vgl.~\cite{Temam1997}, Kapitel~V).  Dies schließt die direkte
Anwendung von Federers Theorem aus.

Ziel dieser Arbeit ist die Einführung einer schwächeren
Regularitätsklasse -- \emph{trajektorienweise log-Lipschitz}-Projektionen
-- die mit fraktaler Geometrie verträglich ist und dennoch für die
Kontrolle der projizierten Trajektorie durch beschränkte Variation (BV)
genügt.  Motiviert durch log-Lipschitz-Regularitätsergebnisse für
dissipative Strömungen
(vgl.~\cite{BaeCannone2016}) beobachten wir, dass die naive globale
log-Lipschitz-Bedingung
\begin{equation}\label{eq:global-LL-intro}
  \|P(x) - P(y)\|
  \;\le\;
  C\,\|x-y\|\,\bigl(1 + \log\tfrac{R}{\|x-y\|}\bigr)
\end{equation}
\emph{nicht} BV von $P \circ u$ für AC-Kurven~$u$ impliziert: Die
Partitionssummen $\sum_i \Delta_i(1+\log(R/\Delta_i))$ divergieren
logarithmisch bei Verfeinerung des Gitters
(Abschnitt~\ref{sec:counterexample}).  Die korrekte Formulierung
ersetzt das Inkrementskalen-Gewicht $\log(R/\|x-y\|)$ durch das
Attraktorabstands-Gewicht $\log(R/d(t))$, wobei
$d(t) = \dist(u(t),\mathcal{A})$.  Dies liefert eine
trajektorienadaptierte Bedingung, deren BV-Summen konvergieren,
wann immer eine Log-Distanz-Integrierbarkeitsbedingung (LDI) erfüllt ist.


\paragraph{Hauptbeiträge.}
Diese Arbeit leistet drei Beiträge: (i)~Wir führen die \emph{trajektorienweise log-Lipschitz}-Bedingung (TLL) und die \emph{Log-Distanz-Integrierbarkeit} (LDI) als die natürliche Regularitätsskala für BV-Projektionen auf fraktale Attraktoren ein, streng zwischen Lipschitz (zu stark) und H\"older (zu schwach). (ii)~Wir beweisen eine BV-Kettenregel (Theorem~\ref{thm:BV-chain}), die zeigt, dass TLL\,+\,LDI für beschränkte-Variations-Selektionsregularität genügt. (iii)~Wir identifizieren hinreichende geometrische Bedingungen (beschränkte Variation der Log-Distanz und Sublevel-Zeitkontrolle), die mit der fraktalen Geometrie dissipativer PDE-Attraktoren verträglich sind.

\begin{remark}[Input-Output-Vertrag]
\label{rem:api-contract}
Das zentrale Ergebnis dieser Arbeit lässt sich als Input-Output-Vertrag formulieren:
\begin{description}
\item[Input:] Eine absolut stetige Trajektorie $u \in AC([0,T]; H)$, ein kompakter Attraktor $\mathcal{A} \subset H$ und eine feste messbare Nächstpunkt-Selektion $P : \mathcal{A}^R \to \mathcal{A}$.
\item[Bedingungen:] (TLL) Die gewählte Abbildung $P$ ist trajektorienweise log-Lipschitz entlang $u$; (LDI) das Log-Distanz-Gewicht $\log(R/d(t))$ ist $\|u'\|$-integrierbar.
\item[Output:] Die zusammengesetzte Selektion $v(t) := P(u(t))$ gehört zu $BV([0,T]; H)$ mit der expliziten Schranke $\mathrm{Var}(v) \leq C_{\mathrm{TLL}} \int_0^T \|u'(t)\| \,\Omega(t)\,dt$.
\end{description}
Dieser Vertrag ist so konzipiert, dass er direkt mit dem Begleitpaper~\cite{FST-NS2026} zusammenwirkt: Der Output (BV-Selektion) ist präzise Annahme~$G'$ dieses Papers.

\smallskip\noindent
\textbf{Assumption~$G'$} (aus~\cite{FST-NS2026}):
\emph{Es existiert eine Nächstpunkt-Selektion
$v(t) = P(u(t)) \in \mathcal{A}$, sodass
$v \in BV([0,T];\,H)$ für jedes $T < T^*$.}

\smallskip\noindent
Kombiniert mit Condition~D und den übrigen Hypothesen aus~\cite{FST-NS2026} (einschließlich der Energieungleichung für den Attraktorabstand) ist Assumption~$G'$ der BV-Selektionsinput in dessen bedingtem Regularitätsargument. Die vorliegende Arbeit beweist nur den Schritt TLL\,+\,LDI $\Rightarrow$ BV. Der Attraktor $\mathcal{A}$ darf fraktale Geometrie besitzen -- TLL und LDI sind mit nicht-ganzzahliger Hausdorff-Dimension verträglich.
\end{remark}


% ==================================================================
\section{Grundlagen und Notation}
% ==================================================================

Durchgehend bezeichnet $H$ einen separablen reellen Hilbertraum mit
Norm $\|\cdot\|$ und innerem Produkt $\langle\cdot,\cdot\rangle$.

\begin{definition}[Dissipativer Semiflow und Attraktor]
\label{def:semiflow}
Ein \emph{dissipativer Semiflow} auf $H$ ist eine Familie
$\{S(t)\}_{t \ge 0}$ stetiger Abbildungen $S(t): H \to H$ mit
$S(0) = \id$, $S(t+s) = S(t) \circ S(s)$, die einen
\emph{kompakten globalen Attraktor}
$\mathcal{A} \subset H$ besitzt: eine kompakte, invariante
($S(t)\mathcal{A} = \mathcal{A}$ für alle $t \ge 0$) Menge, die
jede beschränkte Teilmenge von~$H$ anzieht.
\end{definition}

\begin{definition}[Tubulare Umgebung und Abstand]
\label{def:tube}
Für $R > 0$ ist die \emph{$R$-tubulare Umgebung} von $\mathcal{A}$
\[
  \mathcal{A}^R
  \;:=\;
  \{x \in H : \dist(x, \mathcal{A}) < R\},
\]
wobei $\dist(x, \mathcal{A}) := \inf_{a \in \mathcal{A}} \|x - a\|$.
Für eine Kurve $u: [0,T] \to \mathcal{A}^R$ schreiben wir
$d(t) := \dist(u(t), \mathcal{A})$.
\end{definition}

\begin{definition}[Nächstpunkt-Selektion]
\label{def:selection}
Eine \emph{Nächstpunkt-Selektion} ist eine messbare Abbildung
$P: \mathcal{A}^R \to \mathcal{A}$ mit
$\|x - P(x)\| = \dist(x, \mathcal{A})$ für alle
$x \in \mathcal{A}^R$ und $P|_{\mathcal{A}} = \id$.
Eine solche Selektion existiert nach dem Satz von
Kuratowski--Ryll-Nardzewski, angewandt auf die kompaktwertige
Multifunktion
$x \mapsto \arg\min_{a \in \mathcal{A}} \|x-a\|$.
\end{definition}

\begin{proposition}[Messbare und tie-gebrochene Nächstpunkt-Selektionen]
\label{prop:minimal-selection}
Sei $u \in \AC([0,T];\,H)$ und setze
\[
  \Pi(t) := \{a \in \mathcal{A} : \|u(t)-a\| = d(t)\}.
\]
Dann ist $\Pi(t)$ eine nichtleere, kompaktwertige messbare
Multifunktion und besitzt daher messbare Selektionen
$v(t) \in \Pi(t)$.  Ist ein Borel-Tie-Breaker
$R:\mathcal{A}\to\mathbb{R}$ vorgegeben, der auf $\Pi(t)$ für
f.ü.\ $t$ ein eindeutiges Minimum annimmt (zum Beispiel nach einer
endlich-dimensionalen oder Moreau--Yosida-Regularisierung), dann
definiert
\[
  v_R(t) \in \arg\min_{a\in\Pi(t)} R(a)
\]
auf diesen Zeiten eine kanonisch tie-gebrochene
Nächstpunkt-Selektion.  Auf der verbleibenden mehrdeutigen Menge
kann eine beliebige messbare Selektion unter den Minimierern
gewählt werden.
\end{proposition}

\begin{proof}
Die Nächstpunkt-Menge $\Pi(t) := \{a \in \mathcal{A} :
\|u(t)-a\| = d(t)\}$ ist kompakt und nichtleer für jedes~$t$,
weil $\mathcal{A}$ kompakt ist und $a\mapsto\|u(t)-a\|$ stetig
ist.  Der Graph von $\Pi$ ist messbar, da $u$ messbar ist und die
Distanzfunktion stetig ist.  Der Satz von Kuratowski--Ryll-Nardzewski
liefert daher eine messbare Nächstpunkt-Selektion.  Für einen festen
Borel-Tie-Breaker $R$ ist die eingeschränkte Argmin-Korrespondenz
$t\mapsto\arg\min_{\Pi(t)}R$ wiederum kompaktwertig und messbar;
bei f.ü.\ eindeutigen Minimierern folgt die angegebene
tie-gebrochene Selektion.  Eine f.ü.\ Eindeutigkeit der
Nächstpunkte wird im Allgemeinen nicht behauptet.
\end{proof}

\begin{remark}[Rolle des Tie-Breakings]
Die BV-Kettenregel unten beweist nicht, dass eine bevorzugte
Nächstpunkt-Selektion TLL erfüllt.  Sie setzt vielmehr eine messbare
Selektion $P$ mit TLL voraus und leitet daraus die BV-Aussage ab.
Tie-Breaker wie minimale Enstrophie, lexikographische
Endlich-Moden-Ordnung oder eine regularisierte
Metrische-Ableitungs-Regel sind nützliche Kandidaten zur Konstruktion
einer solchen Selektion; ihre TLL-Regularität ist aber in jedem
konkreten System ein zusätzlicher Satz.
\end{remark}

\begin{remark}[Eindeutigkeit für nichtkonvexe Mengen]\label{rem:uniqueness-caveat}
Für nichtkonvexe $\mathcal{A}$ kann die Nächstpunkt-Menge
$\Pi(t)$ mehr als ein Element enthalten, und weder die strikte
Konvexität der Hilbertnorm noch die Differenzierbarkeit der
eindimensionalen Funktion $t\mapsto d(t)$ erzwingen Eindeutigkeit
des nächsten Punktes in einer nichtkonvexen Menge.  Porositätsresultate
für die mehrdeutige Umgebungsmenge sind daher generische
Hintergrundinformationen, aber kein f.ü.-Eindeutigkeitsargument
entlang jeder absolut stetigen Trajektorie.  Deshalb ist der
Hauptsatz für eine feste Selektion formuliert, die TLL erfüllt.
\end{remark}


% ==================================================================
\section{Warum globale Log-Lipschitz-Stetigkeit versagt}
\label{sec:counterexample}
% ==================================================================

Wir zeigen zunächst, dass die globale log-Lipschitz-Bedingung
\eqref{eq:global-LL-intro}, selbst kombiniert mit
Log-Distanz-Integrierbarkeit, nicht direkt BV von $P \circ u$ liefert.

\begin{proposition}[Versagen der globalen LL-Schranke zur Variationskontrolle]
\label{prop:divergence}
Sei $P: \mathcal{A}^R \to \mathcal{A}$ mit
\eqref{eq:global-LL-intro} und Konstante $C > 0$.  Sei
$u: [0,1] \to \mathcal{A}^R$ Lipschitz mit
$\|u'(t)\| = 1$ f.ü.  Für die gleichmäßige Partition
$t_i = i/n$ ($i = 0, \ldots, n$) erfüllt die aus
\eqref{eq:global-LL-intro} gewonnene obere Schranke
\[
  \sum_{i=0}^{n-1} \|P(u(t_{i+1})) - P(u(t_i))\|
  \;\le\;
  C\,(1 + \log(nR)),
\]
die für $n \to \infty$ divergiert.  Insbesondere liefert die
globale LL-Bedingung keine endliche \emph{a priori}-Schranke
für $\Var(P \circ u)$.
\end{proposition}

\begin{proof}
Jedes Inkrement erfüllt $\Delta_i := \|u(t_{i+1}) - u(t_i)\| = 1/n$.
Nach \eqref{eq:global-LL-intro}:
\[
  \sum_{i=0}^{n-1} \|P(u(t_{i+1})) - P(u(t_i))\|
  \;\le\;
  C \sum_{i=0}^{n-1} \frac{1}{n}\,
  \bigl(1 + \log(nR)\bigr)
  \;=\;
  C\,(1 + \log(nR)).
\]
Für $n \to \infty$ wächst dies wie $C\log n \to \infty$.
Die aus \eqref{eq:global-LL-intro} abgeleitete Schranke divergiert
also mit der Gitterverfeinerung und liefert keine endliche obere
Schranke für $\Var(P \circ u)$.  (Ob die Variation selbst unendlich
ist, hängt von der spezifischen Geometrie von~$\mathcal{A}$ ab;
entscheidend ist, dass die globale LL-Bedingung als \emph{Werkzeug}
zur Herleitung von BV unzureichend ist.)
\end{proof}

\begin{remark}
Die Divergenz ist \emph{logarithmisch}, nicht polynomiell -- viel
langsamer als für H\"older-Abbildungen (wo $\sum \Delta_i^\alpha$
für $\alpha < 1$ konvergiert).  Dies macht globales LL ,,beinahe''
hinreichend für BV, aber die logarithmische Akkumulation über
unendlich viele Partitionspunkte verhindert es.  Die trajektorienweise
log-Lipschitz-Bedingung vermeidet dies, indem sie das Log-Gewicht am
Attraktorabstand $d(t)$ verankert, der nicht von der Partition
abhängt.
\end{remark}


% ==================================================================
\section{Die trajektorienweise log-Lipschitz-Bedingung}
\label{sec:TLL}
% ==================================================================

\begin{definition}[Trajektorienweise log-Lipschitz-Projektion (TLL)]
\label{def:TLL}
Sei $u \in \AC([0,T];\, H)$ mit $u(t) \in \mathcal{A}^R$ für
alle~$t$, und sei $P: \mathcal{A}^R \to \mathcal{A}$ eine
Nächstpunkt-Selektion.  Wir sagen, dass $P$ die
\emph{trajektorienweise log-Lipschitz-Bedingung} (TLL) entlang~$u$
erfüllt, wenn Konstanten $C_{\mathrm{TLL}} > 0$ und $h_0 > 0$
existieren, sodass für alle $t \in [0,T]$ mit $d(t) > 0$ und alle
$h \in \mathbb{R}$ mit $0 < |h| < h_0$ und
$t+h \in [0,T]$
(d.h.\ sowohl Vorwärts- als auch Rückwärts-Inkremente) gilt:
\begin{equation}\label{eq:TLL}
  \|P(u(t+h)) - P(u(t))\|
  \;\le\;
  C_{\mathrm{TLL}}\,\|u(t+h) - u(t)\|\,
  \Omega(t),
\end{equation}
wobei $\Omega(t) := 1 + \log(R/d(t))$ das
\emph{Log-Distanz-Gewicht} ist.

Auf der Menge $\{t : d(t) = 0\}$ (wo $u(t) \in \mathcal{A}$)
gilt $P(u(t)) = u(t)$ durch Idempotenz von~$P$.  Für den
Übergang $d(t) = 0$, $d(t+h) > 0$ gilt
$\|P(u(t+h)) - u(t)\| = \|P(u(t+h)) - P(u(t))\|$.
Da $u(t) \in \mathcal{A}$ ein Kandidat für den nächsten Punkt
von $u(t+h)$ ist, liefert die Minimalität von $P$:
$\|u(t+h) - P(u(t+h))\| \le \|u(t+h) - u(t)\|$.
Zusammen mit der Dreiecksungleichung:
\[
  \|P(u(t+h)) - u(t)\|
  \;\le\;
  \|P(u(t+h)) - u(t+h)\| + \|u(t+h) - u(t)\|
  \;\le\;
  2\,\|u(t+h) - u(t)\|.
\]
Dies ergibt $|v'|(t) \le 2\,|u'|(t)$ auf $\{d = 0\}$.
\end{definition}

\begin{remark}[Konvention \texorpdfstring{$C_{\mathrm{TLL}} \ge 2$}{C_TLL >= 2}]
\label{rem:CTLL-convention}
Der Faktor~$2$ auf $\{d = 0\}$ wird durch die TLL-Schranke
absorbiert, sofern $C_{\mathrm{TLL}} \ge 2$ und $\Omega(t) \ge 1$.
Durchgehend verwenden wir die Konvention
$C_{\mathrm{TLL}} := \max(C_{\mathrm{TLL}},\,2)$;
dies ändert die BV-Schranke \eqref{eq:BV-bound} um höchstens
einen Faktor~$2$ und beeinflusst die Integrierbarkeit
von~\eqref{eq:LDI} nicht.
\end{remark}

\begin{remark}[Geometrische Interpretation von TLL]
\label{rem:TLL-geometry}
Die Bedingung TLL besagt, dass die Nächstpunkt-Projektion,
ausgewertet entlang der Trajektorie~$u$, eine lokale
Lipschitz-Konstante hat, die \emph{höchstens logarithmisch}
divergiert, wenn sich die Trajektorie dem Attraktor nähert:
\[
  \mathrm{Lip}_{\mathrm{loc}}(P \circ u;\, t)
  \;\lesssim\;
  1 + \log\frac{R}{d(t)}.
\]
Im Abstand $d(t) \sim R$ vom Attraktor ist die Projektion
im Wesentlichen Lipschitz (der Log-Term ist $O(1)$).  Für
$d(t) \to 0$ wächst das Log-Gewicht und spiegelt die zunehmende
geometrische Komplexität der Attraktorgrenze auf feinen Skalen
wider.  Die logarithmische Rate ist die \emph{kritische Schwelle}:
Schnelleres Wachstum (z.B. polynomiell) würde LDI generisch
zum Scheitern bringen, während langsameres Wachstum (beschränkt)
Lipschitz-Projektion implizieren würde -- zu stark für fraktale
Geometrie.
\end{remark}

\begin{remark}[TLL als Aubin-Eigenschaft der Projektionskorrespondenz]
\label{rem:TLL-Aubin}
Die TLL-Bedingung lässt eine natürliche Umformulierung in der
Sprache der Variationsanalyse zu.  Definiere die
\emph{Projektionskorrespondenz}
\[
  \Gamma: t \;\mapsto\;
  \bigl\{a \in \mathcal{A} : \|a - u(t)\| = d(t)\bigr\},
\]
die jeder Zeit die Menge der nächsten Attraktorelemente zuordnet.
Die TLL-Bedingung (Definition~\ref{def:TLL}) ist \emph{analog} zur
\emph{Aubin-Eigenschaft} (metrische Regularität) der Korrespondenz
$\Gamma$ mit Rate $1/|\log r|$.  Konkret ist $\Gamma$ metrisch
regulär mit Rate $\phi(r) = 1/|\log r|$, wenn
\[
  \dist\bigl(a,\, \Gamma(t')\bigr)
  \;\le\;
  L\,\phi\bigl(\|u(t') - u(t)\|\bigr)\,\|u(t') - u(t)\|
\]
für $a \in \Gamma(t)$ und $t'$ nahe~$t$ gilt.
TLL und die Aubin-Eigenschaft sind im Allgemeinen \emph{nicht}
äquivalent: TLL gewichtet mit
$\Omega(t) = 1 + \log(R/d(t))$ (dem \emph{Attraktorabstand}),
während die Aubin-Formulierung mit
$1/|\log\|u(t')-u(t)\||$ (der \emph{Inkrementgröße}) gewichtet.
Beide Formulierungen fallen nur in dem Spezialfall zusammen, in dem
$d(t) \sim \|u(t')-u(t)\|$ gilt, also wenn Trajektorien den Attraktor
nicht-tangential erreichen.  Diese Analogie verbindet die
Projektionsregularitätstheorie mit dem Aubin-Kriterium und dem
Rahmenwerk metrischer Regularität von
Rockafellar--Wets~\cite{RockafellarWets1998} und
Dontchev--Rockafellar~\cite{DontchevRockafellar2009}.
Der zentrale Vorteil ist, dass Aubin-Eigenschafts-Kriterien für
mengenwertige Abbildungen gut entwickelt sind und verifizierbare
hinreichende Bedingungen liefern (z.B.\ über die
Mordukhovich-Koderivierte), die auf die Attraktorprojektion
anwendbar sein könnten.
\end{remark}

\begin{remark}[TLL vs.\ globales LL]
\label{rem:TLL-vs-LL}
Die globale log-Lipschitz-Bedingung \eqref{eq:global-LL-intro}
impliziert \emph{nicht} TLL.  Für $h \to 0$ bei festem~$t$:
\[
  \log\frac{R}{\|u(t+h)-u(t)\|}
  \;\approx\;
  \log\frac{R}{h\,|u'(t)|}
  \;=\;
  \log\frac{R}{|u'(t)|} + \log\frac{1}{h}
  \;\to\; \infty,
\]
während $\log(R/d(t))$ fest bleibt.  Die globale Bedingung
erzeugt ein \emph{Inkrementskalen}-Log-Gewicht, das mit der
Partitionsfeinheit divergiert; TLL erzeugt ein
\emph{Abstandsskalen}-Log-Gewicht, das nur von der aktuellen
Geometrie abhängt.  Dies ist der strukturelle Grund, warum TLL
BV liefert, während globales LL dies nicht tut.
\end{remark}


% ==================================================================
\section{Log-Distanz-Integrierbarkeit}
\label{sec:LDI}
% ==================================================================

\begin{definition}[Log-Distanz-Integrierbarkeit (LDI)]
\label{def:LDI}
Im Rahmen von Definition~\ref{def:TLL} sagen wir, dass
\emph{Log-Distanz-Integrierbarkeit} gilt, wenn
\begin{equation}\label{eq:LDI}
  \int_{\{d>0\}} \|u'(t)\|\,\Omega(t)\,\mathrm{d}t
  \;=\;
  \int_{\{d>0\}} \|u'(t)\|\,
  \bigl(1 + \log\tfrac{R}{d(t)}\bigr)\,\mathrm{d}t
  \;<\; \infty.
\end{equation}
Auf der Menge $\{t : d(t) = 0\}$ gilt $v(t) = u(t)$ und der
BV-Beitrag wird durch $\int_{\{d=0\}} \|u'(t)\|\,\mathrm{d}t
< \infty$ kontrolliert (da $u$ AC ist), sodass nur die Region
$\{d > 0\}$ für \eqref{eq:LDI} relevant ist.  Um die
BV-Schranke \eqref{eq:BV-bound} einheitlich auf ganz $[0,T]$
zu formulieren, setzen wir $\Omega(t) := 1$ auf $\{d=0\}$
(konsistent mit der Faktor-$2$-Schranke aus
Bemerkung~\ref{rem:CTLL-convention}).
\end{definition}

\begin{remark}[Minimalität von LDI]
\label{rem:LDI-minimal}
Die Bedingung LDI ist präzise zwischen den bekannten
Regularitätsklassen positioniert:
\begin{itemize}
  \item \textbf{Lipschitz-Projektion} ($\Omega \equiv 1$):
    LDI reduziert sich auf $\int \|u'\| < \infty$, was für
    AC-Kurven automatisch gilt.  Dies ist das Federer-/positive-Reichweite-Regime.
  \item \textbf{H\"older-Projektion} ($\Omega(t) \sim d(t)^{-\beta}$,
    $\beta > 0$): LDI erfordert $\int \|u'\|\,d(t)^{-\beta} < \infty$,
    was generisch in der Nähe des Attraktors scheitert.  Deshalb
    liefern Man\'e-Projektoren (H\"older-Inverse) kein BV.
  \item \textbf{Log-Lipschitz-Projektion} (TLL):
    LDI erfordert $\int \|u'\|\,\log(1/d) < \infty$, den
    Grenzfall -- logarithmische Singularitäten sind unter milden
    dynamischen Bedingungen integrierbar.
\end{itemize}
Das log-Lipschitz/LDI-Paar besetzt somit das \emph{kritische}
Regularitätsfenster für BV-Selektionen auf fraktalen Attraktoren.
\end{remark}


% ==================================================================
\section{Hauptresultat: BV-Kettenregel}
\label{sec:main}
% ==================================================================

\begin{theorem}[BV-Kettenregel für log-Lipschitz-Projektionen]
\label{thm:BV-chain}
Sei $\{S(t)\}_{t \ge 0}$ ein dissipativer Semiflow auf einem
separablen Hilbertraum~$H$ mit kompaktem globalem Attraktor
$\mathcal{A} \subset H$.  Sei $u \in \AC([0,T];\, H)$ mit
$u(t) \in \mathcal{A}^R$ für alle~$t$, und sei
$P: \mathcal{A}^R \to \mathcal{A}$ eine Nächstpunkt-Selektion,
die TLL \eqref{eq:TLL} entlang~$u$ erfüllt.  Setze $v(t) := P(u(t))$.

Wenn LDI \eqref{eq:LDI} gilt, dann ist
$v \in \BV([0,T];\, H)$ mit der expliziten Schranke
\begin{equation}\label{eq:BV-bound}
  \Var_H(v;\,[0,T])
  \;\le\;
  C_{\mathrm{TLL}}
  \int_0^T \|u'(t)\|\,\Omega(t)\,\mathrm{d}t,
\end{equation}
wobei $\Omega(t) := 1 + \log(R/d(t))$ auf $\{d > 0\}$ und
$\Omega(t) := 1$ auf $\{d = 0\}$
(Bemerkung~\ref{rem:CTLL-convention}).
\end{theorem}

\begin{proof}
Der Beweis gliedert sich in drei Schritte.

\medskip
\noindent
\textbf{Schritt 1 (Schranke der metrischen Ableitung).}\;
Für jede stetige Kurve $v: [0,T] \to H$ ist die \emph{obere
metrische Ableitung} bei $t$ definiert als
\[
  |v'|^+(t)
  \;:=\;
  \limsup_{h \to 0} \frac{\|v(t+h) - v(t)\|}{|h|}
  \;\in\; [0,+\infty].
\]
Diese Größe existiert überall (als Wert in $[0,+\infty]$)
und ist Borel-messbar.  Für absolut stetige Kurven ist der
$\limsup$ f.ü.\ ein genuiner Limes
(vgl.~\cite[Theorem~1.1.2]{AmbrosioGigliSavare2008}).
Da $u$ AC ist, existiert $|u'|(t) := \lim_{h\to 0}\|u(t+h)-u(t)\|/|h|$
für f.ü.~$t$.

\smallskip\noindent
\emph{Fall $d(t)>0$:}\;
Division von \eqref{eq:TLL} durch $|h|$ und Bildung des $\limsup$:
\begin{equation}\label{eq:metric-deriv-bound}
  |v'|^+(t)
  \;\le\;
  C_{\mathrm{TLL}}\,|u'|(t)\,\Omega(t)
  \qquad \text{für f.ü.\ } t \text{ mit } d(t) > 0.
\end{equation}
(Die rechte Seite ist ein genuiner Limes; die linke Seite kann
ein $\limsup$ sein, der durch dieselbe Schranke dominiert wird.)

\smallskip\noindent
\emph{Fall $d(t)=0$ (Inneres von $\mathcal{Z}$):}\;
Auf $\mathcal{Z} := \{t \in [0,T] : d(t) = 0\}$ gilt
$u(t) \in \mathcal{A}$ und $v(t) = P(u(t)) = u(t)$.
Da $u(t) \in \mathcal{A}$ ein Kandidat-Minimierer für
$\dist(u(t+h), \mathcal{A})$ ist, liefert die
Dreiecksungleichung
$\|P(u(t+h)) - u(t)\| \le 2\,\|u(t+h) - u(t)\|$
(vgl.\ die Schranke in Definition~\ref{def:TLL}),
also $|v'|^+(t) \le 2\,|u'|(t)$.

\smallskip\noindent
\emph{Randpunkte von $\mathcal{Z}$:}\;
An $t_0 \in \partial\mathcal{Z}$ mit $d(t_0)=0$ können Folgen
$t_n \to t_0$ mit $d(t_n) > 0$ existieren.  Für die Stetigkeit von
$v$ bei $t_0$: $\|v(t_n) - v(t_0)\| = \|P(u(t_n)) - u(t_0)\|
\le d(t_n) + \|u(t_n) - u(t_0)\| \to 0$, da $u$
stetig ist und $d = \dist(\cdot,\mathcal{A}) \circ u$ Stetigkeit
erbt (die Abstandsfunktion ist $1$-Lipschitz).  Die Schranke der metrischen Ableitung
$|v'|^+(t_0) \le 2\,|u'|(t_0)$ bleibt durch dasselbe
Dreiecksungleichungsargument gültig.

\smallskip\noindent
Da $\Omega(t) \ge 1$ und $C_{\mathrm{TLL}} \ge 2$
(Bemerkung~\ref{rem:CTLL-convention}), gilt die
Schranke \eqref{eq:metric-deriv-bound} auf ganz $[0,T]$
(mit $|v'|^+$ statt $|v'|$, wo der Limes möglicherweise
nicht existiert; dies genügt für das BV-Kriterium in Schritt~3).

\medskip
\noindent
\textbf{Schritt 2 (Stetigkeit von~$v$).}\;
Die Nächstpunkt-Projektion auf eine kompakte (aber nichtkonvexe)
Menge ist im Allgemeinen nur oberhalbstetig als mengenwertige
Abbildung und muss nicht überall einwertig sein.  Die
TLL-Bedingung \eqref{eq:TLL} impliziert jedoch direkt die Stetigkeit
von $v = P \circ u$ entlang der Trajektorie: Für $d(t) > 0$ gilt
$\|v(t+h) - v(t)\| \le C_{\mathrm{TLL}}\,\|u(t+h)-u(t)\|\,
\Omega(t) \to 0$ für $h \to 0$; für $d(t) = 0$ gilt
$\|v(t+h)-v(t)\| \le 2\,\|u(t+h)-u(t)\| \to 0$
(einschließlich Randpunkte von $\mathcal{Z}$, vgl.\ Schritt~1).
Also ist $v$ stetig auf $[0,T]$.

\medskip
\noindent
\textbf{Schritt 3 (BV-Schranke).}\;
Wir zeigen $\|v(b) - v(a)\| \le \int_a^b g(t)\,\mathrm{d}t$
für alle $0 \le a < b \le T$, wobei
$g(t) := C_{\mathrm{TLL}}\,|u'|(t)\,\Omega(t) \in L^1$
(nach LDI).  Da $v$ stetig ist (Schritt~2), folgt die
BV-Schranke $\Var(v) \le \int_0^T g$ durch Summation über
jede Partition.

Setze $f(t) := \|v(t) - v(a)\|$.  Dann ist $f: [a,b] \to \mathbb{R}$
stetig und reellwertig,
$f(a) = 0$, und die obere Dini-Ableitung erfüllt
$D^+ f(t) \le |v'|^+(t)$ (da
$f(t+h) - f(t) \le \|v(t+h)-v(t)\|$ nach der umgekehrten
Dreiecksungleichung, und Division durch $h > 0$ mit
$\limsup$ ergibt $D^+ f \le |v'|^+$).
Nach Schritt~1 gilt $|v'|^+(t) \le g(t)$ f.ü.  Ein klassisches Resultat über Dini-Ableitungen
(vgl.\ \cite[Kapitel~VII, \S\S\,9--10]{Saks1937}; siehe
auch \cite[Kapitel~3]{AmbrosioFuscoPallara2000}) besagt:
Für eine stetige Funktion $f:[a,b] \to \mathbb{R}$ mit
$D^+ f(t) \le g(t)$ f.ü.\ und $g \in L^1$ gilt die
Ungleichung $f(b) - f(a) \le \int_a^b g(t)\,\mathrm{d}t$.
Daher
\[
  \|v(b) - v(a)\| \;=\; f(b)
  \;\le\; \int_a^b g(t)\,\mathrm{d}t
  \;\le\; \int_0^T g(t)\,\mathrm{d}t.
\]
Summation über jede Partition $0 = t_0 < \cdots < t_n = T$:
\[
  \sum_i \|v(t_{i+1}) - v(t_i)\|
  \;\le\;
  \sum_i \int_{t_i}^{t_{i+1}} g(t)\,\mathrm{d}t
  \;=\;
  \int_0^T g(t)\,\mathrm{d}t
  \;=\;
  C_{\mathrm{TLL}} \int_0^T |u'|(t)\,\Omega(t)\,\mathrm{d}t,
\]
was \eqref{eq:BV-bound} ergibt.

\medskip
\noindent
Die Schranke \eqref{eq:BV-bound} folgt aus den Schritten~1--3.
\end{proof}

\begin{remark}[Direkte Verifikation über Partitionssummen]
\label{rem:partition-sum}
Für eine feste Partition mit Feinheit $< h_0$ lässt sich die
Endlichkeit der Partitionssumme direkt aus TLL ablesen, ohne die
Dini-Ableitungs-Maschinerie: Auf jedem Teilintervall liefert TLL
$\|v(t_{i+1}) - v(t_i)\|
\le C_{\mathrm{TLL}}\,\|u(t_{i+1}) - u(t_i)\|\,\Omega(t_i)$
(oder die Faktor-$2$-Schranke auf $\{d = 0\}$, absorbiert durch
$C_{\mathrm{TLL}} \ge 2$).  Die resultierende Partitionssumme
$\sum_i C_{\mathrm{TLL}}\,\Omega(t_i)\int_{t_i}^{t_{i+1}}
\|u'\|\,\mathrm{d}t$ ist für jede feste Partition endlich
(da $\Omega$ auf jeder kompakten Menge disjunkt von $\{d=0\}$
beschränkt ist und der $\{d=0\}$-Beitrag durch $2\int\|u'\|$
kontrolliert wird).  Der Übergang zum Supremum über \emph{alle}
Partitionen -- also zur Definition von
$\Var(v)$ -- erfordert jedoch das Dini-Ableitungs-Argument
von Schritt~3, das das punktweise $\Omega(t_i)$-Gewicht
zum integrierten $\Omega(t)$-Gewicht verschärft.
\end{remark}

\begin{remark}[Banachraum-Allgemeinheit]
\label{rem:banach}
Der Beweis von Satz~\ref{thm:BV-chain} verwendet nur die Norm,
Kompaktheit von~$\mathcal{A}$ und messbare Selektion -- kein
inneres Produkt und keine Hilbertraum-Struktur.  Die
BV-Kettenregel gilt daher wortwortlich in jedem separablen
Banachraum.  Wir beschränken uns nur wegen der PDE-Anwendungen
in Abschnitt~\ref{sec:applications} auf Hilberträume.
\end{remark}


% ==================================================================
\section{Hinreichende Bedingungen für LDI}
\label{sec:sufficient}
% ==================================================================

Die Bedingung LDI \eqref{eq:LDI} ist eine einzelne
Integrierbarkeitsanforderung.  Wir geben zwei unabhängige
hinreichende Bedingungen an, die jeweils aus dynamischen
Informationen über den Semiflow verifizierbar sind.

% ------------------------------------------------------------------
\subsection{Variante A: Log-Distanz hat beschränkte Variation}
% ------------------------------------------------------------------

\begin{proposition}[BV-Log-Distanz impliziert LDI]
\label{prop:variant-A}
Angenommen $d(t) > 0$ für f.ü.\ $t \in [0,T]$ und
$\log(R/d(t)) \in \BV([0,T])$.  Dann gilt LDI.
\end{proposition}

\begin{proof}
Schreibe $\ell(t) := 1 + \log(R/d(t))$.  Da $u' \in L^1(0,T;\,H)$
(weil $u$ AC ist) und $\ell \in \BV([0,T]) \subset
L^\infty([0,T])$, ist das Produkt
$\|u'(t)\|\,\ell(t)$ lokal integrierbar:
\[
  \int_0^T \|u'(t)\|\,\ell(t)\,\mathrm{d}t
  \;\le\;
  \|\ell\|_{L^\infty(0,T)}\,\int_0^T \|u'(t)\|\,\mathrm{d}t
  \;<\; \infty.  \qedhere
\]
\end{proof}

\begin{remark}[Wann ist die Log-Distanz BV?]
\label{rem:log-BV}
Die Bedingung $\log(R/d) \in \BV$ erfordert, dass die Trajektorie
nicht unendlich oft zwischen verschiedenen Abstandsskalen zum
Attraktor oszilliert.  Formal schließt sie unendliche
,,Zickzack-Frequenz'' in $d(t)$ aus.  Für monoton anziehende
Trajektorien ($d(t)$ schließlich fallend) ist dies automatisch
erfüllt.  Für Trajektorien mit transienten Oszillationen folgt
BV von $\log d$, wenn die Gesamtzahl der Skalenkreuzungen
$\{t : d(t) = \varepsilon\}$ für jedes
$\varepsilon > 0$ endlich ist -- eine generische Eigenschaft
dissipativer Semiflows mit isolierten Gleichgewichten.
\end{remark}

% ------------------------------------------------------------------
\subsection{Variante B: Sublevel-Zeitmaß-Kontrolle}
% ------------------------------------------------------------------

\begin{definition}[Sublevel-Zeit-Kontrolle (STC)]
\label{def:STC}
Die Trajektorie~$u$ erfüllt \emph{Sublevel-Zeit-Kontrolle} mit
Exponent $\alpha > 0$, wenn ein $C_* > 0$ existiert mit
\begin{equation}\label{eq:STC}
  \bigl|\{t \in [0,T] : d(t) \le \varepsilon\}\bigr|
  \;\le\;
  C_*\,\varepsilon^\alpha
  \qquad \text{für alle } \varepsilon \in (0, R].
\end{equation}
\end{definition}

\begin{remark}[STC impliziert \texorpdfstring{$|\{d=0\}| = 0$}{Mass von d=0 ist null}]
\label{rem:STC-d0}
Der Grenzübergang $\varepsilon \to 0$ in \eqref{eq:STC} zeigt
$|\{d = 0\}| = 0$.  STC ist daher nur auf Trajektorien anwendbar,
die keine positive Zeit auf dem Attraktor verbringen.  Dies ist der
generische Fall für nicht-stationäre Lösungen, die sich
$\mathcal{A}$ asymptotisch nähern.  Für Trajektorien, die den
Attraktor in endlicher Zeit erreichen (z.B.\ Konvergenz zu einem
Gleichgewicht mit $u(t) \in \mathcal{A}$ für $t \ge t_1$), ist
Variante~A (BV-Log-Distanz) die geeignete hinreichende Bedingung.
\end{remark}

\begin{proposition}[STC impliziert LDI]
\label{prop:variant-B}
Falls STC \eqref{eq:STC} mit einem $\alpha > 0$ gilt und
$u' \in L^2(0,T;\,H)$, dann ist
$\log(R/d) \in L^p(0,T)$ für jedes $p < \infty$, und LDI gilt.
\end{proposition}

\begin{proof}
\textbf{Schritt 1 (Tonelli-Darstellung).}\;
Da $u(t) \in \mathcal{A}^R$, gilt $d(t) \le R$ für alle~$t$.
Nach Bemerkung~\ref{rem:STC-d0} impliziert STC
$|\{d=0\}|=0$. Alle logarithmischen Integrale unten werden daher über
$\{d>0\}$ genommen; Trajektorien mit Attraktorkontakt positiven Maßes
gehören zu Variante~A, nicht zu STC. Auf $\{t : d(t) > 0\}$ ergibt die
Schichtzerlegungsidentität
$\log(R/d(t)) = \int_{d(t)}^{R} s^{-1}\,\mathrm{d}s$ mit Tonelli:
\begin{equation}\label{eq:tonelli}
  \int_{\{d>0\}} \|u'(t)\|\,\log\frac{R}{d(t)}\,\mathrm{d}t
  \;=\;
  \int_0^R \frac{1}{s}
  \Bigl(\int_{\{0 < d(t) \le s\}} \|u'(t)\|\,\mathrm{d}t\Bigr)
  \,\mathrm{d}s.
\end{equation}

\medskip
\noindent
\textbf{Schritt 2 (H\"older auf dem inneren Integral).}\;
Nach der H\"olderschen Ungleichung:
\[
  \int_{\{0<d(t)\le s\}} \|u'(t)\|\,\mathrm{d}t
  \;\le\;
  \Bigl(\int_0^T \|u'(t)\|^2\,\mathrm{d}t\Bigr)^{1/2}
  \cdot
  \bigl|\{0 < d \le s\}\bigr|^{1/2}
  \;\le\;
  \|u'\|_{L^2}\,C_*^{1/2}\,s^{\alpha/2},
\]
wobei der letzte Schritt $|\{0<d\le s\}| \le |\{d\le s\}|
\le C_*\,s^\alpha$ aus STC verwendet.

\medskip
\noindent
\textbf{Schritt 3 (Integrierbarkeit nahe $s=0$).}\;
Einsetzen in \eqref{eq:tonelli}:
\[
  \int_0^R \frac{1}{s}\,\|u'\|_{L^2}\,C_*^{1/2}\,s^{\alpha/2}
  \,\mathrm{d}s
  \;=\;
  \|u'\|_{L^2}\,C_*^{1/2}
  \int_0^R s^{\alpha/2 - 1}\,\mathrm{d}s
  \;=\;
  \|u'\|_{L^2}\,C_*^{1/2}\,
  \frac{2}{\alpha}\,R^{\alpha/2}
  \;<\; \infty.
\]
Das Integral konvergiert, weil $\alpha > 0$, sodass
$s^{\alpha/2-1}$ nahe $s=0$ integrierbar ist.

\medskip
\noindent
\textbf{Schritt 4 ($L^p$-Schranken).}\;
Dasselbe Tonelli-Argument mit H\"older-Exponenten $p$ und $q$
($1/p + 1/q = 1$) ergibt
$\|\log(R/d)\|_{L^p} \le C(p,\alpha,R,C_*)$ für jedes
$p < \infty$, sodass LDI aus jeder endlichen $L^q$-Schranke
für~$\|u'\|$ folgt.
\end{proof}

\begin{remark}[Quellen von STC in dissipativen Strömungen]
\label{rem:STC-sources}
Die Sublevel-Zeit-Kontrolle \eqref{eq:STC} ist der genuine neue
analytische Input.  Drei Mechanismen können sie erzeugen:

\begin{enumerate}[label=(\roman*)]
  \item \textbf{Untere Distanzhülle auf endlichen Fenstern:}\;
    Eine einseitige Attraktionsabschätzung $d(t)\le C_0 e^{-\lambda t}$
    kontrolliert $\log(R/d(t))$ nicht von oben und impliziert für sich
    genommen keine STC.  Benötigt wird eine nicht-superexponentielle
    untere Hülle, etwa $d(t)\ge c_T e^{-\lambda_T t}$ auf einem festen
    Zeitfenster $[0,T]$.  Dann kann $\{d\le\varepsilon\}$ erst nach der
    Schwelle $(\log(c_T/\varepsilon))/\lambda_T$ auftreten; dies liefert
    STC-Konstanten auf dem festen Fenster (abhängig von
    $T,c_T,\lambda_T$) und direkter die logarithmische Schranke aus
    Lemma~\ref{lem:stc-tracking}.
  \item \textbf{Lyapunov-Dissipation:}\;
    Falls ein Lyapunov-Funktional $E$ die Bedingung
    $-\dot{E} \ge \Psi(d)$ mit nicht zu flachem $\Psi$ nahe $d=0$
    erfüllt (z.B.\ $\Psi(d) \ge c\,d^\beta$ für ein $\beta$),
    dann werden die Sublevel-Zeiten durch das gesamte Energiebudget
    $E(0) - \inf E$ kontrolliert.
  \item \textbf{Keine schnellen Oszillationen:}\;
    Energetische Einschränkungen (endliches Dissipationsintegral)
    verhindern beliebig schnelle Hin-und-Her-Bewegungen nahe dem
    Attraktor und begrenzen die Anzahl der Exkursionen in
    $\{d \le \varepsilon\}$.
\end{enumerate}
Jeder Mechanismus liefert STC mit verschiedenen Exponenten~$\alpha$.
Die STC-Bedingung ist physikalisch plausibel für alle
Standardklassen dissipativer PDEs.
\end{remark}

\begin{lemma}[Untere Exponentialhülle impliziert Log-Integrierbarkeit]
\label{lem:stc-tracking}
Seien $0<c\le R$ und $\lambda\ge 0$, und auf $[0,T]$ gelte
\[
  0<d(t)\le R,
  \qquad
  d(t)\ge c\,e^{-\lambda t}.
\]
Dann ist $\log(R/d)\in L^\infty(0,T)$, und
\[
  \int_0^T \log\frac{R}{d(t)}\,\mathrm{d}t
  \;\le\;
  T\log\frac{R}{c} + \frac{\lambda T^2}{2}.
\]
Falls zusätzlich $u'\in L^1(0,T;\,H)$ gilt, dann gilt LDI auf $[0,T]$.
\end{lemma}

\begin{proof}
Die untere Hülle liefert
\[
  \log\frac{R}{d(t)}
  \;\le\;
  \log\frac{R}{c\,e^{-\lambda t}}
  \;=\;
  \log\frac{R}{c}+\lambda t.
\]
Daher
\[
  \int_0^T \log\frac{R}{d(t)}\,\mathrm{d}t
  \;\le\;
  \int_0^T \bigl(\log(R/c)+\lambda t\bigr)\,\mathrm{d}t
  \;=\;
  T\log\frac{R}{c}+\frac{\lambda T^2}{2}.
\]
Für die LDI-Bedingung \eqref{eq:LDI}:
\[
  \int_0^T \|u'(t)\|\,\log\frac{R}{d(t)}\,\mathrm{d}t
  \;\le\;
  \Bigl(\log\frac{R}{c}+\lambda T\Bigr)
  \|u'\|_{L^1(0,T;H)}
  \;<\; \infty
\]
weil der logarithmische Faktor auf dem endlichen Fenster beschränkt ist.
\end{proof}

\begin{remark}[Exponentialhüllen und STC]
\label{rem:exp-tracking-strongest}
Lemma~\ref{lem:stc-tracking} ist bewusst mit einer unteren Hülle
formuliert.  Eine reine Attraktionsabschätzung
$d(t)\le C_0e^{-\lambda t}$ gibt nur eine obere Schranke für die
Distanz zum Attraktor; sie verhindert nicht, dass $d(t)$ viel kleiner
wird, und kann daher $\log(R/d(t))$ nicht von oben beschränken.
Exponentielle Attraktion ist also nützliche Evidenz, aber LDI benötigt
entweder direkte STC, BV-Kontrolle der Log-Distanz oder eine untere
Hülle beziehungsweise eine Nicht-superexponentialitätsabschätzung.
\end{remark}

\begin{remark}[STC als realistischer Kandidat für Navier--Stokes]
\label{rem:stc-realistic}
Von den drei hinreichenden Bedingungen für LDI (Variante~A: BV-Log-Distanz; Variante~B: Sublevel-Zeitkontrolle; Variante~C: untere exponentielle Distanzhülle) ist die \emph{Sublevel-Zeitkontrolle} (STC) der realistischste Kandidat für 3D-Navier--Stokes:
\begin{enumerate}
\item \emph{Exponentielle Attraktion} ($d(t) \leq C_0 e^{-\lambda t}$) ist das einfache dynamische Regime fern von Blow-up, aber für sich genommen kein LDI-Beweis. LDI verwendet entweder direkte STC oder eine untere Distanzhülle wie in Lemma~\ref{lem:stc-tracking}.
\item \emph{STC} ($|\{d(t) \leq \varepsilon\}| \leq C_* \varepsilon^\alpha$) ist robust: Es erfordert lediglich, dass die Trajektorie auf keiner Skala zu lange nahe dem Attraktor verweilt. Zwei PDE-Mechanismen machen STC für Navier--Stokes plausibel:
\begin{itemize}
\item \emph{Dissipative Zeitskalen-Separation:} Nahe dem Attraktor ist die Niedermoddynamik langsam (kontrolliert durch den Attraktor), während die Hochmoddynamik schnell ist (viskose Dämpfung). Diese Separation begrenzt die Verweilzeit auf jeder Distanzskala.
\item \emph{Endlich viele bestimmende Moden:} Die Anzahl der bestimmenden Moden $N_d \sim G^{2/3}$ (wobei $G$ die Grashof-Zahl ist) begrenzt die effektive Dimension der Dynamik nahe $\mathcal{A}$ und verhindert pathologische Akkumulation auf jeder Distanzskala.
\end{itemize}
\item \emph{BV-Log-Distanz} (Variante~A) erfordert monotone Annäherung an den Attraktor, was generisch durch oszillierende Lösungen verletzt wird.
\end{enumerate}
\end{remark}


% ==================================================================
\section{Anwendungen auf dissipative PDEs}
\label{sec:applications}
% ==================================================================

Wir illustrieren, wie der abstrakte Rahmen auf konkrete
PDE-Situationen anwendbar ist.  In jedem Fall besitzt der Semiflow
einen kompakten globalen Attraktor, und die Frage ist, ob TLL und
LDI gelten. Für dreidimensionale Navier-Stokes muss dieser bedingte
starke Semiflow-Rahmen vom unbedingten Leray-Hopf- beziehungsweise
Trajektorienattraktor-Rahmen unterschieden werden.

% ------------------------------------------------------------------
\subsection{3D-Navier-Stokes auf \texorpdfstring{$\mathbb{T}^3$}{T3}}
% ------------------------------------------------------------------

\begin{example}[Navier-Stokes-Gleichungen]
\label{ex:NS}
Sei $H = L^2(\mathbb{T}^3;\,\mathbb{R}^3)$ (divergenzfrei) und
$f$ glatt. Die folgende Anwendung ist bedingt darauf zu lesen, dass
man in einem starken einwertigen Semiflow oder in einem vorab
gewählten Trajektorienattraktor-Rahmen arbeitet, in dem dieselben
Kompaktheits- und Selektionsobjekte fixiert sind. In der unbedingten
Leray-Hopf-Theorie verhindert die Nicht-Eindeutigkeit, dass man einfach einen
Temam-artigen globalen starken Attraktor als gewöhnliches
Semiflow-Objekt aufruft. Im bedingten starken Attraktor-Ast liefert
die Standard-Attraktortheorie:
\begin{itemize}
  \item $\mathcal{A}$ ist kompakt in $H^1$ (also auch in $H$).
  \item Jedes $a \in \mathcal{A}$ ist $C^\infty$.
  \item Exponentielles Tracking gilt \emph{bedingt auf globale
    Regularität}: Für starke Lösungen, die auf $[0,\infty)$
    existieren, gilt $d(t) \le C_0\,e^{-\lambda t}$ für $t \ge t_0$.
    (Für schwache Leray-Hopf-Lösungen ist nur algebraische
    Attraktion in $L^2$ unbedingt bekannt.)  Diese Attraktionsabschätzung
    ist für sich genommen kein LDI-Beweis; sie muss mit STC,
    BV-Log-Distanz-Kontrolle oder einer unteren Hülle kombiniert werden.
\end{itemize}

In Regimen, in denen die endlichdimensionalen Attraktorabschätzungen
gelten, hat der Attraktor endliche fraktale Dimension
$d_f(\mathcal{A}) \le c_1\,G^{2/3}(1+\log G)^{1/3}$
(wobei $G$ die Grashof-Zahl ist;
vgl.~\cite{ConstantinFoiasTemam1985}), hat aber generisch
Reichweite~$0$ -- was Federers Theorem ausschließt.

\medskip\noindent
\textbf{Status von TLL:}\; Offen.  Die Squeezing-Eigenschaft des
Navier-Stokes-Semiflows (vgl.~\cite{FoiasManleyRosaTemam2001},
Kapitel~14) liefert exponentielle Kontraktion hochfrequenter
Differenzen auf dem Attraktor. Kombiniert mit
Man\'e-Projektor-Abschätzungen deutet dies auf eine mögliche
endlichmodige Route zu TLL hin, kontrolliert aber für sich genommen
nicht die Nächstpunkt-Abbildung. Ein rigoroser Beweis erfordert
quantitative Kontrolle der Projektionsstabilität im
Hochmoden-Komplement $Q_N H$ und des gewählten Nächstpunkt-Zweigs.

\medskip\noindent
\textbf{Status von LDI:}\; Offen für das kritische 3D-Problem.
LDI folgt auf festen Fenstern aus STC oder aus dem Unterhüllen-Kriterium
von Lemma~\ref{lem:stc-tracking}, aber nicht aus der einseitigen
Attraktionsabschätzung $d(t)\le C_0e^{-\lambda t}$ allein.  Die
kritische Frage betrifft das Blow-up-Regime $T \nearrow T^*$, wo
$\|u'\| \to \infty$ und die Distanzgeometrie degenerieren kann.
Wenn TLL und LDI in einem nicht-zirkulären gewählten Rahmen
bewiesen sind, liefert die BV-Schranke \eqref{eq:BV-bound}
Annahme~G$'$ aus \cite{FST-NS2026}; das Begleitpaper nutzt diese
Annahme dann als Teil seines bedingten Attraktorabstands-Arguments.

\medskip\noindent
\textbf{Verbindung zum Log-Dirichlet-Quotienten:}\;
Das Begleitpaper~\cite{FST-NS2026} führt den
\emph{Log-Dirichlet-Quotienten}
$Q_{\mathrm{LD}}(t) = \|\nabla(u-v)\|^2 / \|u-v\|^2$ ein und zeigt,
dass dessen Integrierbarkeit die Gronwall-Abschätzung von
exponentiellem zu algebraischem Wachstum transformiert. Die
LDI-Bedingung der vorliegenden Arbeit liefert den
Selektionsvariations-Input für diese Route: Die BV-Kettenregel
kontrolliert die relevante Log-Distanz-Variation. Eine echte
$Q_{\mathrm{LD}}$-Schranke benötigt weiterhin die Vergleichbarkeits-
und Energiehypothesen des Begleitframeworks.
\end{example}

% ------------------------------------------------------------------
\subsection{Reaktions-Diffusions-Systeme}
% ------------------------------------------------------------------

\begin{example}[Reaktions-Diffusion auf beschränkten Gebieten]
\label{ex:RD}
Betrachte $\partial_t u = D\Delta u + f(u)$ auf einem beschränkten
Gebiet $\Omega \subset \mathbb{R}^n$ mit Dirichlet- oder
Neumann-Randbedingungen, wobei $D > 0$ eine Diffusionsmatrix ist und
$f$ eine Dissipativitätsbedingung
$\langle f(u), u \rangle \le C - c|u|^2$ erfüllt.  Der Semiflow
besitzt einen kompakten globalen Attraktor
$\mathcal{A} \subset L^2(\Omega)$ mit endlicher fraktaler
Dimension~\cite{Temam1997}.

Für skalare Gleichungen ($u \in \mathbb{R}$) in einer Raumdimension
mit Gradienten-Nichtlinearität $f(u) = -V'(u)$ besteht der
Attraktor aus Gleichgewichten und heteroklinen Orbits.  In Regimen,
in denen der Inertialmannigfaltigkeits-Zweig hinreichend regulär ist
(etwa prox-regulär oder $C^{1,1}$ in einer Röhrenumgebung), ist die
Nächstpunkt-Projektion auf diesem Zweig Lipschitz; dann folgen TLL
und LDI im Sinne von Theorem~\ref{thm:squeezing-tll}.
Inertialmannigfaltigkeitsresultate vom Mallet-Paret--Sell-Typ liefern
die natürliche Testumgebung~\cite{MalletParetSell1988}.

Für Systeme in höheren Dimensionen kann der Attraktor
nicht-ganzzahlige fraktale Dimension und Reichweite null haben.
Der log-Lipschitz-Rahmen bietet dann einen Weg zu BV-Selektionen,
den Federers Theorem nicht bereitstellt.  Die Spektrallückenbedingung
für Inertialmannigfaltigkeiten ist für viele
Reaktions-Diffusions-Systeme erfüllt~\cite{EdenFoiasNicolaenkoTemam1994},
was die Attraktorgeometrie regulärer macht als im
Navier-Stokes-Fall.
\end{example}

% ------------------------------------------------------------------
\subsection{Gedämpfte Wellengleichungen}
% ------------------------------------------------------------------

\begin{example}[Gedämpfte semilineare Wellengleichung]
\label{ex:wave}
Betrachte $\partial_{tt} u + \gamma\,\partial_t u = \Delta u + f(u)$
auf $\Omega \subset \mathbb{R}^3$ mit $\gamma > 0$ und subkritischer
Nichtlinearität.  Der Semiflow auf $H^1_0 \times L^2$ besitzt einen
kompakten globalen Attraktor mit endlicher fraktaler Dimension
(vgl.~\cite{BabinVishik1992}).  Die exponentielle Tracking-Eigenschaft
gilt mit Rate $\lambda = \gamma/2$ (Spektrallücke des gedämpften
Wellenoperators; vgl.~\cite{BabinVishik1992}, Kapitel~V).  Für LDI
ist diese Attraktionsabschätzung weiterhin nur ein dynamischer Input;
sie muss wie oben mit STC oder einer unteren Distanzhülle kombiniert
werden.

Die Attraktorgeometrie liegt zwischen Reaktions-Diffusion
(wo Inertialmannigfaltigkeiten oft existieren) und Navier-Stokes
(wo sie nicht existieren).  Der log-Lipschitz-Rahmen bietet einen
\emph{einheitlichen} Zugang, der auf alle drei Situationen
anwendbar ist, mit TLL als einzigem geometrischen Input.
\end{example}


% ==================================================================
\section{Diskussion}
\label{sec:discussion}
% ==================================================================

\subsection{Die neue mathematische Schnittstelle}

Der zentrale Beitrag dieser Arbeit ist die Identifikation der
\emph{kritischen Regularitätsschnittstelle} für BV-Selektionen
auf fraktalen Attraktoren.  Die Schnittstelle wird durch eine
einzige Frage charakterisiert:

\begin{center}
\fbox{\parbox{0.85\textwidth}{%
\textbf{Kernfrage.}\;
Erfüllt für einen dissipativen Semiflow mit kompaktem Attraktor
$\mathcal{A} \subset H$ die Nächstpunkt-Projektion die
trajektorienweise log-Lipschitz-Kontrolle \eqref{eq:TLL}?
}}
\end{center}

\noindent
Dies ist eine Frage über die \emph{Regularität des Attraktors
aus Sicht der Trajektorie}, nicht über die globale Geometrie
von~$\mathcal{A}$.  Die Unterscheidung ist entscheidend: Globale
geometrische Bedingungen (positive Reichweite, Inertialmannigfaltigkeit)
erzwingen Struktur auf dem \emph{gesamten} Attraktor, während TLL
Regularität nur in Richtung der dynamischen Annäherung erfordert.

\subsection{Bezug zu bestehenden Regularitätstheorien}

\begin{center}
\resizebox{\textwidth}{!}{%
\begin{tabular}{lllll}
\hline
\textbf{Bedingung} & \textbf{Projektion} & \textbf{Komposition}
  & \textbf{Fraktal?} & \textbf{Quelle} \\
\hline
Positive Reichweite $> 0$
  & Lipschitz-1
  & AC $\Rightarrow$ AC
  & Nein
  & Federer~\cite{Federer1959} \\
Prox-regulär
  & lokal Lipschitz
  & AC $\Rightarrow$ AC (lokal)
  & Nein
  & Poliquin et al.~\cite{PoliquinRockafellarThibault2000} \\
Hausdorff-BV-Mengenabb.
  & BV-Selektion existiert
  & ---
  & Ja
  & Chistyakov~\cite{Chistyakov2004} \\
Man\'e-Projektor
  & H\"older${}^{-1}$
  & AC $\Rightarrow$ H\"older
  & Ja
  & Man\'e~\cite{Mane1981} \\
\textbf{TLL + LDI}
  & \textbf{log-Lipschitz}
  & \textbf{AC $\Rightarrow$ BV}
  & \textbf{Ja}
  & \textbf{Diese Arbeit} \\
\hline
\end{tabular}%
}
\end{center}

\noindent
Der TLL+LDI-Rahmen ist der einzige Eintrag, der gleichzeitig
(i)~fraktale Geometrie berücksichtigt (nicht-ganzzahlige Dimension,
Reichweite~$0$), (ii)~BV-Regularität liefert (hinreichend für
Gronwall-artige Argumente) und (iii)~als Trajektorien-Bedingung
formuliert ist (unabhängig von der globalen Attraktorstruktur).

\subsection{Einschränkungen}

\begin{itemize}
  \item \textbf{Selektionsabhängigkeit.}\;
    Die TLL-Bedingung ist für eine \emph{gegebene} messbare
    Nächstpunkt-Selektion~$P$ formuliert.  Auf nichtkonvexen
    kompakten Mengen ist die Nächstpunkt-Abbildung generisch
    mehrwertig, und verschiedene Selektionen können unterschiedliche
    Regularität haben.  Der Rahmen erfordert die Existenz
    \emph{mindestens einer} Selektion, die TLL erfüllt; er
    behauptet nicht, dass jede messbare Selektion dies tut.
    In der Praxis sind vorregistrierte tie-gebrochene Selektionen
    (z.B.\ minimale Enstrophie, lexikographische Endlich-Moden-Ordnung
    oder eine regularisierte Metrische-Ableitungs-Regel) die natürlichen
    Kandidaten.
  \item \textbf{TLL ist offen für 3D-NS.}\;
    Die Squeezing-Eigenschaft und Man\'e-Projektor-Abschätzungen
    deuten auf log-Lipschitz-Stabilität der Projektion hin, aber ein
    rigoroser Beweis erfordert quantitative Schranken für das
    Hochmoden-Projektionskomplement, die derzeit nicht verfügbar sind.
  \item \textbf{LDI ist nur unter Blow-up nicht-trivial.}\;
    Für glatte Lösungen auf kompakten Zeitintervallen gilt LDI
    automatisch.  Die Bedingung wird erst im Grenzfall
    $T \nearrow T^*$ (Blow-up-Zeit) kritisch, wo $\|u'\| \to \infty$.
  \item \textbf{BV, nicht AC.}\;
    Der Rahmen liefert BV-Regularität, die für
    Stieltjes-Integral-basierte Gronwall-Argumente genügt, aber nicht
    für punktweise Differentialungleichungen.  Ein Upgrade von BV auf
    AC würde Lipschitz-Projektion (positive Reichweite) erfordern --
    genau die Bedingung, die wir zu vermeiden suchen.
  \item \textbf{BV-Schranke hängt von $\|u'\|$ ab.}\;
    Die Schranke \eqref{eq:BV-bound} lautet
    $\Var(v) \le C_{\mathrm{TLL}} \int_0^T \|u'\|\,\Omega\,dt$.
    Für Navier--Stokes hängt $\|u'\|$ von der Lösungsnorm ab,
    und nahe einem potentiellen Blow-up ($T \nearrow T^*$) gilt
    $\|u'\| \to \infty$. Der Gronwall-Abschluss erfordert, dass die
    BV-Schranke ,,absorbiert'' werden kann, bevor $\Var(v)$ divergiert.
    Die Details dieser Absorption hängen vom
    Begleitpaper~\cite{FST-NS2026} ab und liegen außerhalb des
    Rahmens dieser Arbeit.
\end{itemize}


\subsection*{Zentrales offenes Problem}
\begin{quote}
\textbf{Selektionsregularitätsproblem.}
In einem starken Navier--Stokes-Semiflow oder in einer vorab gewählten Trajektorienattraktor-Formulierung für Leray--Hopf-Lösungen: Erfüllt die Nächstpunkt-Projektion oder die gewählte Nächstpunkt-Korrespondenz $P_\mathcal{A}$ die TLL-Bedingung entlang $u(t)$? Äquivalent: Erlaubt die Attraktorgeometrie nicht-zirkuläre BV-reguläre Selektionen?

Eine bejahende Antwort, kombiniert mit der BV-Kettenregel dieser Arbeit, würde Assumption~$G'$ für das Begleitframework~\cite{FST-NS2026} liefern. Jeder Schluss auf globale Regularität hängt weiterhin zusätzlich von Condition~D und den übrigen Hypothesen dieses bedingten Frameworks ab.
\end{quote}


% ==================================================================
\section{Schlussfolgerung}
% ==================================================================

Wir haben einen Regularitätsrahmen für Nächstpunkt-Projektionen
auf kompakte Attraktoren eingeführt, der die kritische Lücke
zwischen H\"older (unzureichend für BV) und Lipschitz (unvereinbar
mit fraktaler Geometrie) besetzt.  Die trajektorienweise
log-Lipschitz-Bedingung (TLL) kontrolliert den Projektionsmodul auf
der Attraktorabstandsskala, und die
Log-Distanz-Integrierbarkeitsbedingung (LDI) stellt sicher, dass die
resultierenden BV-Summen konvergieren.

Der Rahmen reduziert bedingte Regularitätskriterien für dissipative
PDEs -- einschließlich starker oder vorab gewählter
dreidimensionaler Navier-Stokes-Formulierungen -- auf eine
geometrisch-dynamische Frage: Hat eine feste Nächstpunkt-Selektion
höchstens logarithmisches Blow-up entlang dissipativer Trajektorien?
Diese Frage ersetzt nicht die PDE-spezifischen Kompaktheits-,
Selektions- und Grenzübergangshypothesen; sie isoliert den
geometrischen Teil dieser Last.


% ==================================================================
\section{Zeitgemittelte Selektion und BV-Regularität (spekulativ)}
\label{sec:time-averaged}
% ==================================================================

\emph{Dieser Abschnitt ist spekulativ: Er skizziert einen alternativen Ansatz,
dessen Kernschritte offene Vermutungen bleiben. Die Resultate der
Abschnitte~\ref{sec:main}--\ref{sec:applications} hängen nicht von
diesem Material ab.}

\medskip

Die bevorzugte Transferarchitektur ist inzwischen der endlichmodige
Skeleton-Grenzübergang mit vorregistrierten regularisierten Selektoren
(Abschnitt~\ref{sec:zookeeper-ldi-transfer}). Die folgende
Zeitmittelungskonstruktion bleibt nur als ergänzende Selektionsheuristik
erhalten.

\medskip

Die zentrale Obstruktion zur Schließung der Regularitätslücke ist
die BV-Kontrolle der Nächstpunkt-Selektion
$v(t) \in \arg\min_{a \in \mathcal{A}} \|u(t) - a\|$.  Wir schlagen
vor, die punktweise Selektion durch eine \emph{zeitgemittelte}
Selektion zu ersetzen, die nach zusätzlichem geometrischem Input
besseres BV-Verhalten haben kann.

\begin{definition}[Moreau--Yosida-regularisierte Selektion]
\label{def:my-selection}
Für $\varepsilon > 0$ definiere die \emph{zeitgemittelte
Minimierermenge}
\[
  \Pi_\varepsilon(t)
  \;:=\;
  \arg\min_{a \in \mathcal{A}}
  \int_{t-\varepsilon}^{t+\varepsilon}
  \|u(s) - a\|^2 \, ds.
\]
Jedes gewählte Element $v_\varepsilon(t)\in\Pi_\varepsilon(t)$ muss
durch einen vorregistrierten messbaren Tie-Breaker ausgewählt werden.
\end{definition}

\begin{conjecture}[Eigenschaften von $v_\varepsilon$]
\label{conj:bv-regularity}
Unter den Standardannahmen (endlich-dimensionaler Attraktor,
Energiedissipationsschranke) ist eine vorregistrierte regularisierte
Selektion $v_\varepsilon(t)\in\Pi_\varepsilon(t)$ ein Kandidat für die
folgenden Eigenschaften:
\begin{enumerate}[label=(\roman*)]
\item \emph{Eindeutigkeit:}  $v_\varepsilon(t)$ ist f.ü.\ eindeutig
  (strikte Konvexität des zeitintegrierten Kostenfunktionals auf~$H$;
  für die Minimierung über das nicht-konvexe Kompaktum~$\mathcal{A}$
  erfordert Eindeutigkeit ein separates Argument über
  Hausdorff-Stabilität der Minimierer-Menge, das offen bleibt).
\item \emph{BV-Schranke bei festem \(\varepsilon\):}
  $\Var(v_\varepsilon;\,[0,T])
   \;\le\;
   C_\varepsilon
   \int_0^T \|\partial_t u(s)\|\,ds
   \;<\; \infty$,
  für einen Zweig, für den die Sensitivitätsabschätzung unten gilt,
  wobei $C_\varepsilon$ typischerweise für $\varepsilon \to 0$
  divergiert. Eine \emph{uniforme} BV-Schranke (unabhängig
  von~$\varepsilon$) ist eine zusätzliche geometrische Hypothese, keine
  Konsequenz der Definition; siehe Bemerkung~\ref{rem:epsilon-limit}.
\item \emph{TLL-Vererbung:}  Falls~$u$ TLL mit Konstante
  $C_{\mathrm{TLL}}$ erfüllt, wird erwartet, dass $v_\varepsilon$ TLL mit
  Konstante $C_{\mathrm{TLL}} + O(\varepsilon)$ erfüllt. Ein rigoroser
  Beweis erfordert Kontrolle der Diskrepanz zwischen punktweisem und
  gemitteltem Abstand, was offen bleibt.
\item \emph{Gronwall-Kompatibilität:}  Die Gronwall-Abschätzungen
  des Hauptresultats (Satz~\ref{thm:BV-chain}) bleiben voraussichtlich für
  $v_\varepsilon$ gültig mit einem zusätzlichen Fehlerterm
  $O(\varepsilon\,\|\partial_t u\|_{L^2})$. Eine rigorose Formulierung
  erfordert die Spezifikation des relativen Entropiefunktionals~$H$ und
  der Konstanten $\alpha, \beta, \gamma$; dies wird auf zukünftige Arbeit verschoben.
\end{enumerate}
\end{conjecture}

\begin{remark}[Partielle Heuristik]
\emph{(i)}\; Die Abbildung
$a \mapsto \int_{t-\varepsilon}^{t+\varepsilon} \|u(s)-a\|^2\,ds$
ist strikt konvex auf~$H$. Die Minimierung erfolgt jedoch über das
kompakte, im Allgemeinen \emph{nicht-konvexe} $\mathcal{A}$.
Eindeutigkeit des Minimierers über einer nicht-konvexen Menge erfordert
ein separates Argument (z.B.\ über Hausdorff-Stabilität der
Minimierer-Menge oder ein generisches Transversalitätsargument).
Dieser Schritt bleibt für allgemeine Attraktoren offen; er gilt,
wenn~$\mathcal{A}$ eine Lipschitz-Mannigfaltigkeit ist.

\emph{(ii)}\; Falls ein gewählter Zweig eindeutig und unter
Zeitverschiebungen lokal stabil ist, legt die übliche
First-Variation-Sensitivität die Abschätzung
\[
  \|v_\varepsilon(t+h) - v_\varepsilon(t)\|
  \;\le\;
  \frac{1}{2\varepsilon}
  \int_{t-\varepsilon}^{t+\varepsilon}
  \|u(s+h) - u(s)\|\,ds
  \;\le\;
  \frac{h}{2\varepsilon}
  \int_{t-\varepsilon}^{t+\varepsilon}
  \|\partial_t u(s)\|\,ds.
\]
Integration über $t \in [0,T]$ und Fubini würde eine
BV-Schranke bei festem \(\varepsilon\) mit einer Konstante der
Größenordnung $\varepsilon^{-1}$ ergeben. Dies ist kein Beweis für
beliebige nicht-konvexe Attraktoren; es notiert nur die Abschätzung,
die für einen stabil gewählten Zweig zu rechtfertigen wäre.

\emph{(iii)}\; Die TLL-Konstante verschlechtert sich um das
Zeitmittelungsfenster:
\[
  |\log d_\varepsilon(t+h) - \log d_\varepsilon(t)|
  \;\le\;
  C_{\mathrm{TLL}}\,|\log|h||
  \;+\;
  C\,\varepsilon/|h|,
\]
wobei der zweite Term aus der Diskrepanz zwischen punktweisem und
gemitteltem Abstand entsteht.

\emph{(iv)}\; Die Gronwall-Abschätzung nimmt einen zusätzlichen
Forcingterm aus der Zeitmittelung auf:
\[
  \frac{d}{dt} H(u \mid v_\varepsilon)
  \;\le\;
  -\alpha\,H + \beta + \gamma\,\varepsilon,
\]
wobei $\gamma$ von $\|\partial_t u\|$ abhängt, aber nicht von~$H$.
Das Standard-Gronwall-Lemma ist anwendbar; der zusätzliche
$\gamma\varepsilon$-Term verschwindet im Limes
$\varepsilon \to 0$.
\end{remark}

\begin{remark}[Grenzübergang \texorpdfstring{$\varepsilon \to 0$}{epsilon to 0}]
\label{rem:epsilon-limit}
Die zentrale Frage ist, ob die BV-Schranke den Grenzübergang
$\varepsilon \to 0$ überlebt:
\begin{itemize}
\item Falls\; $\sup_\varepsilon \Var(v_\varepsilon;\,[0,T]) < \infty$,
  konvergiert eine Teilfolge $v_{\varepsilon_k}$ in~$L^1$ gegen eine
  BV-Selektion~$v_0$, die die \emph{kanonische} Nächstpunkt-Selektion
  ist.
\item Eine gleichmäßige BV-Schranke würde einer neuen geometrischen
  Bedingung entsprechen, die man ,,zeitgemittelte Reichweite'' nennen
  kann. Der Ausdruck ist hier nur ein Name für die fehlende
  Stabilitätshypothese; es wird keine Äquivalenz zu einem bekannten
  Reichweitenbegriff behauptet.
\end{itemize}
Dies ist die präzise Formulierung der ,,BV-Regularitätslücke'' und
verbindet sich mit:
\begin{itemize}
\item \emph{Moreau--Yosida-Regularisierung} in der konvexen Analysis
  (die zeitgemittelten Kosten sind die Moreau-Enveloppe),
\item \emph{Viskoser Selektion} in Gradientensystemen (die Zeitmittelung
  wirkt als künstliche Viskosität),
\item \emph{Thermodynamischer Selektion} im Dark-Energy-Begleitpaper
  (die IR-Mittelung von~$\Phi[\rho]$).
\end{itemize}
\end{remark}


% ==================================================================
\section{Post-Zookeeper-Transfer: Vom Proof of Life zu gedämpften 3D-Navier--Stokes}
\label{sec:zookeeper-ldi-transfer}

Die Zookeeper-Normalform macht dieses Paper zum natürlichen Testträger
des Navier--Stokes-Asts.  Hier sind Defekt, Gap und Selektionsmechanismus
bereits explizit:

In der aktualisierten CORE-Lesart heißt das: RFEP v1.8 importiert eine
Normalform-Checkliste aus dem Zookeeper-Abschluss von Riemann
$C^{(2)}$/MS2.  Das ist noch kein Beweis der dreidimensionalen
Navier--Stokes-Regularität; der Transfer von gedämpften zu klassischen
Systemen bleibt die astspezifische Pflicht.
\begin{description}[style=nextline,leftmargin=0pt,labelsep=0pt]
  \item[Operator/Form.] Die Abstand-zum-Attraktor-Abbildung und ihre
  messbaren nearest-point-Selektionen.
  \item[Defekt.] Das Log-Distanz-Integral, die Projektionsvariation und
  der BV-Defekt.
  \item[Approximation.] Lorenz $\to$ Kuramoto--Sivashinsky $\to$
  Reaktions-Diffusion $\to$ gedämpfte 3D-Navier--Stokes $\to$
  klassische 3D-Navier--Stokes.
  \item[Gap.] TLL, Squeezing, untere Distanzhüllen und Damping-Gaps.
  \item[Grenzübergang.] Von endlichdimensionalen Attraktoren zu
  PDE-Attraktoren und von gedämpften\slash hyperviskosen Systemen
  zur klassischen Gleichung mit $\beta=1$.
\end{description}

Das unmittelbare nächste Ziel ist daher nicht das volle
Millenniumsproblem.  Es ist ein Transfersatz für gedämpfte
3D-Navier--Stokes-Gleichungen: In einem Parameterbereich, in dem die
Literatur zu dynamischen Systemen einen globalen Attraktor und verbesserte
Regularität liefert, soll Dämpfung zusammen mit unabhängig
kontrollierter Projektionszweig-Geometrie und LDI eine BV-Selektion
implizieren.  Die bestehenden Lorenz- und Kuramoto--Sivashinsky-Tests
werden dadurch zu den ersten zwei Stufen einer größeren Leiter, nicht zu
isolierter numerischer Evidenz.

Ein nicht-zirkulärer Transfer sollte daher als
Skeleton-Kompaktheitsproblem formuliert werden. Für ein regularisiertes
System setze \(A_{\delta,N}=P_N\mathcal{A}_\delta\) und definiere einen
vorregistrierten Selektor
\[
v_{\delta,N,\eta}(t)\in\arg\min_{a\in A_{\delta,N}}
\frac12\|P_Nu_\delta(t)-a\|^2+\eta R_N(a),
\]
wobei \(R_N\) vor Beobachtung des Erfolgs fixiert wird, etwa als
endlichmodige lexikographische Ordnung, minimale Enstrophie oder
minimale metrische Ableitung. Der benötigte Satz lautet nicht:
gedämpfter Erfolg impliziert klassischen Erfolg. Er lautet:
Skeleton-Konvergenz, uniforme Skeleton-TLL, uniforme
Tie-/Sublevel-Zeitkontrolle und uniforme BV-Schranken implizieren eine
BV-Grenzselektion für das klassische Problem.

\section{Offene Richtungen}
\label{sec:open-directions}
% ==================================================================

Der TLL+LDI-Rahmen eröffnet mehrere konkrete Forschungsrichtungen,
die thematisch gruppiert sind.

% ------------------------------------------------------------------
\subsection{Pilotsysteme zur Verifikation von TLL und LDI}
% ------------------------------------------------------------------

\begin{remark}[Pilotsysteme zur Verifikation von TLL und LDI]
\label{rem:pilot-verification}
Drei Kandidatensysteme zur Pilotverifikation der TLL- und LDI-Bedingungen:
\begin{enumerate}
\item \emph{Reaktions-Diffusion (2D):} Glatte Nichtlinearität, gut verstandene Attraktoren. TLL und LDI gelten erwartungsgemäß aufgrund der parabolischen Glättung.
\item \emph{Hyperviskose NSE ($(-\Delta)^\beta$, $\beta \geq 5/4$):} Kandidatisches Brückensystem, in dem Inertialmannigfaltigkeits- und Glättungsabschätzungen die zusätzliche Röhrenzweig-Kontrolle für TLL liefern könnten. Dies ist eine Testroute, keine automatische TLL-Garantie.
\item \emph{Kuramoto--Sivashinsky (1D):} Deterministisches Chaos ohne Singularitäten. Der Attraktor ist endlich-dimensional und glatt, was TLL+LDI durch direkte Berechnung verifizierbar macht.
\end{enumerate}
\end{remark}

\begin{remark}[Proof of Life: TLL und LDI auf dem Lorenz-Attraktor]
\label{rem:tll-ldi-lorenz}
Wir verifizieren TLL und LDI numerisch auf dem Lorenz-Attraktor ($\sigma = 10$, $\rho = 28$, $\beta = 8/3$), einem 3D-chaotischen System mit fraktaler Attraktordimension $\approx 2.06$. Der Attraktor wird durch $50\,000$ Trajektorienpunkte approximiert; eine Testtrajektorie wird mit Störung $(2, 2, 2)$ von einem Attraktorelement aus initialisiert und für $T = 100$ Zeiteinheiten integriert.

\emph{Ergebnisse} (die vier beweisrelevanten Diagnosen bestanden; die
Exponentialzeile ist eine separate Annäherungsraten-Diagnose):

\begin{center}
\small
\resizebox{\textwidth}{!}{%
\begin{tabular}{l|l|l}
\textbf{Test} & \textbf{Ergebnis} & \textbf{Detail} \\
\hline
TLL (Definition~\ref{def:TLL}) & PASS & $C_{\mathrm{TLL}}(95\%) = 0{,}19$, max $= 1{,}94$ \\
LDI (Definition~\ref{def:LDI}) & PASS & $\int \|u'\| \Omega\,\mathrm{d}t = 6{,}47 \times 10^4$ (endlich) \\
BV-Kettenregel (Theorem~\ref{thm:BV-chain}) & PASS & $\mathrm{Var}(v) / (C_{\mathrm{TLL}} \cdot \mathrm{LDI}) = 0{,}52$ \\
STC (Definition~\ref{def:STC}) & PASS & $\alpha = 1{,}11 > 0$ \\
Exponentielle Annäherungsrate & DIAGNOSTIK & $\lambda = 0{,}064 > 0$; nicht als LDI-Beweis verwendet
\end{tabular}%
}
\end{center}

Wesentliche Beobachtungen: (i) Die TLL-Konstante ist bemerkenswert klein ($C_{\mathrm{TLL}} \approx 0{,}19$ am 95. Perzentil), was bestätigt, dass die Nächste-Punkt-Projektion entlang chaotischer Trajektorien gutartig ist; (ii) das LDI/L$^1$-Verhältnis beträgt $6{,}86$, was zeigt, dass das Log-Distanz-Gewicht das Integral um weniger als eine Größenordnung vergrößert; (iii) die BV-Kettenregelschranke wird mit Spielraum erfüllt (Verhältnis $0{,}52 < 1$); (iv) Sublevel-Time-Control gilt mit Exponent $\alpha \approx 1{,}1$, konsistent mit der erwarteten Potenzgesetz-Annäherung an den Attraktor.

\emph{Limitierungen des numerischen Tests:}
Die Attraktor-Approximation verwendet $50\,000$ Trajektorienpunkte,
was $d(t)$ relativ zum wahren Attraktor systematisch überschätzt
(was TLL-Konstanten kleiner erscheinen lässt). Zudem handelt es
sich um einen einzelnen Test bei den klassischen Lorenz-Parametern
($\sigma=10$, $\rho=28$, $\beta=8/3$) in $\mathbb{R}^3$ -- einer
endlichdimensionalen ODE, keiner unendlichdimensionalen PDE.
Ein systematischer Parameterscan über $\rho$ (einschließlich
Near-Onset-Werten $\rho \approx 24$) würde die Evidenz
verstärken. Trotz dieser Vorbehalte bestätigt der Test, dass die
TLL+LDI-Bedingungen nicht leer sind: Sie sind in mindestens einem
konkreten chaotischen dynamischen System erfüllt.
\end{remark}

\begin{remark}[Proof of Life~II: TLL und LDI auf dem Kuramoto--Sivashinsky-Attraktor]
\label{rem:tll-ldi-ks}
Wir erweitern die numerische Validierung auf die Kuramoto--Sivashinsky (KS)-Gleichung
$u_t + u\,u_x + u_{xx} + u_{xxxx} = 0$ auf $[0, L]$ mit periodischen Randbedingungen
($L = 32\pi$), eine 1D dissipative PDE mit raumzeitlichem Chaos und
endlichdimensionalem Attraktor. Die pseudospektrale Diskretisierung verwendet $N = 128$
Fourier-Moden; der Attraktor wird durch $5\,001$ Punkte im
$M = 32$-dimensionalen Fourier-Koeffizientenraum approximiert.

\emph{Ergebnisse} (primärer Lauf, $N = 128$):

\begin{center}
\small
\resizebox{\textwidth}{!}{%
\begin{tabular}{l|l|l}
\textbf{Test} & \textbf{Ergebnis} & \textbf{Detail} \\
\hline
TLL (Definition~\ref{def:TLL}) & PASS & $C_{\mathrm{TLL}}(95\%) = 9{,}0$, max $= 58{,}8$ \\
LDI (Definition~\ref{def:LDI}) & PASS & $\int \|u'\| \Omega\,\mathrm{d}t = 5576$ (endlich), LDI/$L^1$ $= 1{,}86$ \\
BV-Kettenregel (Theorem~\ref{thm:BV-chain}) & PASS & $\mathrm{Var}(v) / (C_{\mathrm{TLL}} \cdot \mathrm{LDI}) = 0{,}059$ \\
STC (Definition~\ref{def:STC}) & PASS & $\alpha = 0{,}91 > 0$ \\
Exponentielles Tracking & N/A & $\lambda = -0{,}074$ (chaotisches Wandern, keine monotone Annäherung)
\end{tabular}%
}
\end{center}

\emph{Gitterverfeinerung} ($N = 32, 64, 128, 256$; $N=32$ verwendet $L = 22$):

\begin{center}
\small
\resizebox{\textwidth}{!}{%
\begin{tabular}{r|r|r|r|r|r|c}
$N$ & $M$ & $C_{\mathrm{TLL}}(95\%)$ & LDI/$L^1$ & BV-Verh. & STC $\alpha$ & Alle \\
\hline
32  & 16 & 1{,}81 & 2{,}58 & 0{,}145 & 1{,}85 & PASS \\
64  & 32 & 8{,}75 & 1{,}41 & 0{,}083 & 1{,}20 & PASS \\
128 & 32 & 5{,}10 & 2{,}02 & 0{,}078 & 1{,}11 & PASS \\
256 & 32 & 2{,}66 & 2{,}77 & 0{,}099 & 0{,}79 & PASS
\end{tabular}%
}
\end{center}

Wesentliche Beobachtungen: (i)~Alle vier Diagnosen bestehen bei jeder Auflösung -- das
TLL/LDI-Framework ist \emph{kein} Artefakt endlicher Diskretisierung.
(ii)~Das BV-Verhältnis ist gleichmäßig klein ($< 0{,}15$), was zeigt, dass die
Kettenregelschranke weit vom Optimum entfernt ist.
(iii)~Anders als das Lorenz-System ($\mathbb{R}^3$) ist KS eine \emph{PDE}
mit unendlichdimensionalem Phasenraum, trunkiert auf $N$ Fourier-Moden.
Die Attraktordimension ($\sim 8$--$12$ für $L = 32\pi$) liegt deutlich über
dem Lorenz-Wert ($\sim 2{,}06$), was dies zu einem wesentlich härteren Test macht.
(iv)~Exponentielles Tracking schlägt fehl (erwartet: Die Testtrajektorie wandert
\emph{auf} dem Attraktor, nicht zu ihm hin).
Skript: \texttt{compute\_tll\_ldi\_ks.py}; Ergebnisse: \texttt{ks\_tll\_ldi\_results.json}.
\end{remark}

\begin{remark}[Ginzburg--Landau als TLL/LDI-Brückenkandidat]
\label{rem:ginzburg-landau}
Die komplexe Ginzburg--Landau-Gleichung $\partial_t u = (1+i\alpha)\Delta u + u - (1+i\beta)|u|^2 u$ besitzt für bestimmte Parameterbereiche eine endlich-dimensionale Inertialmannigfaltigkeit. Vollständige Verifikation von TLL+LDI auf diesem System, gefolgt von einem Lifting-Argument auf 2D-NS (via Parameterfortsetzung in der Nichtlinearität), würde einen konkreten Testweg für das NS-LDI-Transferprogramm liefern. Dies ist nicht als direkter Weg zum vollen klassischen 3D-NS-Problem ohne die oben isolierten uniformen Projektionszweig-Abschätzungen zu lesen.
\end{remark}

\begin{theorem}[Konditionales Squeezing-zu-TLL-Kriterium]\label{thm:squeezing-tll}
Sei $u_t + Au + B(u,u) = f$ eine dissipative PDE auf einem Hilbertraum $H$
mit kompaktem globalem Attraktor $\mathcal{A}$ endlicher fraktaler Dimension $d_F$.
Die Squeezing-Eigenschaft gelte: Es existieren $N \in \mathbb{N}$,
$\theta \in (0,1)$ und $T > 0$, sodass für alle $u_0, v_0$ in der
absorbierenden Kugel
\[
\begin{aligned}
  \|Q_N(S(T)u_0 - S(T)v_0)\|
  &\leq \theta\,\|u_0 - v_0\|,\\
  \text{oder}\qquad
  \|Q_N(S(T)u_0 - S(T)v_0)\|
  &\leq \|P_N(S(T)u_0 - S(T)v_0)\|.
\end{aligned}
\]
wobei $P_N$ auf die ersten $N$ Eigenmoden von $A$ projiziert.
Zusätzlich sei entlang der betrachteten Trajektorie ein einwertiger
Nächstpunkt-Zweig $\pi_{\mathcal{A}}$ auf einer Röhre
$\{x:\dist(x,\mathcal{A})<r_0\}$ gegeben, die das betrachtete
Trajektoriensegment enthält, und dieser Zweig sei dort mit Konstante
$L_\pi$ Lipschitz.  Dann gilt die Trajektorien-Log-Lipschitz-Bedingung
(TLL, Definition~\ref{def:TLL}) auf diesem Segment mit
\[
  C_{\mathrm{TLL}} \leq \max\{L_\pi,2\}.
\]
Sind die Röhrenzweig-Konstanten durch die Squeezing-Daten kontrolliert,
kann man dies schematisch als
$C_{\mathrm{TLL}}\le C(N,d_F,\theta,r_0,L_\pi)$ schreiben.
\end{theorem}

\begin{proof}
Das Argument verläuft in zwei Schritten.

\emph{Schritt~1 (Lipschitz-Graph).}
Die Squeezing-Eigenschaft impliziert, dass der globale Attraktor $\mathcal{A}$
in einem Lipschitz-Graphen über dem endlich-dimensionalen Unterraum $P_N H$ liegt.
Dies ist das Man\'e-Projektionstheorem angewandt auf die Squeezing-Dichotomie:
Da die Hochmoden-Komponente $Q_N$ relativ zur Niedrigmoden-Komponente $P_N$
kontrahiert, ist die Einschränkung $P_N|_{\mathcal{A}}$ injektiv mit
Lipschitz-Inverser \cite{FoiasTemam1979,Mane1981}.
Konkret existiert $L = L(N, d_F) < \infty$, sodass
$\|Q_N(u - v)\| \leq L\,\|P_N(u - v)\|$ für alle $u, v \in \mathcal{A}$.

\emph{Schritt~2 (Der fehlende Projektionsregularitätsinput).}
Der vorige Schritt allein impliziert \emph{nicht}, dass die
Nächstpunkt-Projektion auf $\mathcal{A}$ Lipschitz ist: Ein allgemeiner
Lipschitz-Graph muss weder positive Reichweite besitzen noch
prox-regulär sein.  Genau dies ist die eigentliche Lücke.  Unter der
zusätzlichen Annahme eines Lipschitz-Röhrenzweigs gilt jedoch für
$t,t+h$ im betrachteten Segment
\[
  \|\pi_{\mathcal{A}}(u(t+h))-\pi_{\mathcal{A}}(u(t))\|
  \le L_\pi\,\|u(t+h)-u(t)\|.
\]
Da $\Omega(t)\ge 1$ ist, ist dies gerade TLL mit
$C_{\mathrm{TLL}}\le L_\pi$ außerhalb von $\{d=0\}$; die Konvention
$C_{\mathrm{TLL}}\ge 2$ behandelt die Menge $\{d=0\}$ wie in
Bemerkung~\ref{rem:CTLL-convention}.  Squeezing identifiziert also
die endlich-dimensionale Graphroute, aber prox-reguläre oder sonst
Lipschitz-kontrollierte Nächstpunkt-Zweige sind der zusätzliche
Baustein, der daraus TLL macht.
\end{proof}

\begin{remark}[Hyperviskose Navier--Stokes als konditionaler TLL/LDI-Pilot]
\label{rem:hyperviscous}
Für die hyperviskose NS-Gleichung $\partial_t u + (-\Delta)^\beta u + (u\cdot\nabla)u = -\nabla p + f$
mit $\beta \geq 5/4$ auf $\mathbb{T}^3$ machen die bekannte
Glättung und die Literatur zu Inertialmannigfaltigkeiten dieses System
zu einem starken Piloten für das Kriterium oben.  Wenn der relevante
Inertialmannigfaltigkeits-Zweig uniforme Nächstpunkt-Röhrenkontrolle
und unabhängige LDI-Kontrolle liefert (etwa über STC oder eine untere
Distanzhülle), erhält man die konditionale Kette
\[
\begin{aligned}
  \text{IM/prox-regulärer Zweig}
  &\Rightarrow \text{Squeezing + Zweigkontrolle}\\
  &\Rightarrow \text{TLL}
  \Rightarrow \text{TLL + LDI}
  \Rightarrow \text{BV-Selektion (G')}\\
  &\Rightarrow \text{konditionale Regularität im Begleitpaper}.
\end{aligned}
\]
Dies ist als proof-grade Testroute zu lesen, nicht als unbedingter
Abschluss des klassischen 3D-Navier--Stokes-Problems.  Für klassische
NS ($\beta = 1$) bleiben der uniforme Nächstpunkt-Zweig/TLL-Input und
die entsprechende LDI-Kontrolle offen.
\end{remark}

% ------------------------------------------------------------------
\subsection{Mathematische Werkzeuge für TLL und LDI}
% ------------------------------------------------------------------

\begin{remark}[Brückenheuristik: Squeezing zu TLL]
\label{rem:bridge-squeezing}
Die Squeezing-Eigenschaft (Foias--Temam) besagt, dass hohe Moden schneller kontrahieren als niedrige Moden divergieren. Kombiniert mit bestimmenden Moden kann sie Abschätzungen der Form
\[
  d(u(t),\mathcal{A})
  \le C e^{-\gamma_N t}d(u(0),\mathcal{A})
      + C'\sup_{s\le t}\|P_N u(s)-P_N v(s)\|
\]
für eine geeignete Vergleichstrajektorie $v\in\mathcal{A}$ liefern.
Solche Abschätzungen sind nützlicher Input für Log-Distanz-Schranken,
beweisen aber nicht für sich allein Nächstpunkt-TLL.  Sie müssen mit
der Röhren-Projektionszweig-Regularität aus
Theorem~\ref{thm:squeezing-tll} kombiniert werden.
\end{remark}

\begin{remark}[Verdopplungs- und Assouad-Struktur]
\label{rem:doubling-assouad}
Assouad-artige Kontrolle der Attraktortrajektorie kann einen Weg zu
TLL liefern: Wenn die Trajektorie $\{u(t):t\in[0,T]\}$ hinreichend
dünne metrische Struktur besitzt, können approximative Tangentialkegel
logarithmische Öffnungsschranken der Ordnung $C/|\log r|$ haben.  Dies
ist ein plausibler hinreichender Mechanismus für TLL, aber keine
Äquivalenz in dieser Arbeit.  Neuere Strukturresultate über Mengen
positiver Reichweite~\cite{Lytchak2024} liefern zusätzliche Werkzeuge
zum Verständnis der Geometrie von Attraktorumgebungen.
\end{remark}

\begin{remark}[LDI als Schwanzintegral-Äquivalenz]
\label{rem:ldi-tail}
Die LDI-Bedingung
\[
  \int_{\{d>0\}} \|u'(t)\|\,\log\frac{R}{d(t)}\,\mathrm{d}t < \infty
\]
ist äquivalent zur Schwanzintegral-Bedingung
\[
  \int_0^R \frac{1}{s}
  \left(
    \int_{\{0 < d \leq s\}} \|u'(t)\|\,\mathrm{d}t
  \right)\mathrm{d}s < \infty.
\]
Sie setzt die gewichtete Integrierbarkeit von $\log(R/d)$ mit der
,,Geschwindigkeit'' der Trajektorie nahe dem Attraktor in Beziehung
und ist für Energieabschätzungen besser geeignet.
\end{remark}

\begin{remark}[Klärung zu Proposition~\ref{prop:variant-B}]
\label{rem:prop-variant-B-clarification}
Die BV $\to$ $L^\infty$-Einbettung in Proposition~\ref{prop:variant-B} erfordert die zusätzliche Bedingung $d(t) > 0$ f.ü. Die Trajektorie darf den Attraktor nicht auf einer Menge positiven Maßes berühren. Dies ist eine Kein-Kontakt-positiven-Maßes-Annahme, kein automatischer Leray-Hopf-Fakt. Für starke analytische Lösungen kann sie unter zusätzlichen Hypothesen gerechtfertigt werden; auf Intervallen, auf denen die Trajektorie auf dem Attraktor liegt, ist stattdessen Variante~A zu verwenden.
\end{remark}

% ------------------------------------------------------------------
\subsection{Stochastische Erweiterungen und fortgeschrittene Werkzeuge}
% ------------------------------------------------------------------

\begin{remark}[Stochastische Stabilisierung]
\label{rem:stochastic-stabilisation}
Hinzufügen von Feller-artigem Rauschen $\sigma\,\mathrm{d}W$ zu den Navier--Stokes-Gleichungen kann helfen, TLL-artige Transversalität ,,fast sicher'' zu erzwingen: Das Rauschen kann dauerhaft tangentiale Annäherung an den Attraktor reduzieren, und die Feller-Eigenschaft liefert Stetigkeit der Übergangshalbgruppe. Dies ist eine heuristische Stabilisierungsroute, kein Satz, dass stochastischer Antrieb automatisch einen TLL-Nächstpunkt-Zweig liefert.
\end{remark}

\begin{remark}[Hochmoden-Stabilitätsvermutung]
\label{rem:high-mode-stability}
Ein möglicher Weg zu TLL im NS-Kontext ist Hochmoden-Stabilität: Wenn die hohen Fourier-Moden $Q_N u = (I - P_N)u$ uniforme Abfallabschätzungen erfüllen und der dadurch induzierte Nächstpunkt-Zweig im entsprechenden Niedrigmoden-Skeleton kontrolliert ist, kann TLL aus bestimmenden-Moden-Abschätzungen plus dem Röhrenzweig-Kriterium folgen. Hochmoden-Abfall allein genügt nicht, weil er Distanz zu einem Niedrigmoden-Skeleton kontrolliert, nicht die Regularität der Nächstpunkt-Abbildung. Neuere scharfe Dimensionsabschätzungen für NS-artige Attraktoren~\cite{IlyinKalantarovZelik2025} liefern quantitativen Input für diese Route.
\end{remark}

\begin{remark}[Perspektive der Variationsanalysis]
\label{rem:variational-analysis}
Die Theorie der metrischen Regularität und Variationsanalysis (Rockafellar--Wets~\cite{RockafellarWets1998}, Dontchev--Rockafellar~\cite{DontchevRockafellar2009}) bietet eine natürliche Heimat für die TLL- und LDI-Bedingungen: TLL steht in Beziehung zur Aubin-Eigenschaft der Attraktorprojektion, und LDI ist eine Calmness-Bedingung für die Distanzfunktion. Das vollständige Werkzeugset der Variationsanalysis -- einschließlich Ekelands Variationsprinzip, des Mordukhovich-Kriteriums und der Attouch--Th\'era-Dualität -- ist anwendbar und in diesem Kontext wenig erforscht.
\end{remark}

% ------------------------------------------------------------------
\subsection{Weitere NS-spezifische Forschungsrichtungen}
% ------------------------------------------------------------------

\begin{remark}[Weitere Forschungsrichtungen]
\label{rem:further-ns-logdist}
\begin{enumerate}
\item \emph{Selektionsregularität über Trajektorienattraktoren:} Formulierung der minimalen metrischen Ableitungsselektion im Chepyzhov--Vishik-Trajektorienattraktor-Rahmen, wo Nichteindeutigkeit durch Betrachtung ganzer Trajektorien statt Anfangsdaten behandelt wird.
\item \emph{Quantitative Tracking-Lemmata:} Herleitung von Maßabschätzungen $|\{t : d(t) \leq \varepsilon\}| \leq C \varepsilon^\alpha$ über Energie-/Absorbierende-Kugel-Argumente und ,,Kein-klebriger-Kontakt''-Bedingungen.
\item \emph{Ekelands Variationsprinzip:} Nutzung von Ekelands Prinzip zur Konstruktion von ,,Fast-Minimierern'' mit kontrollierter Variation, wodurch die Notwendigkeit positiver Reichweite von $\mathcal{A}$ umgangen wird.
\item \emph{Mengenwertige BV-Selektion:} Bestimmung, welche dynamischen Regularitätsbedingungen an den Fluss beschränkte Hausdorff-Variation der Minimierermenge $\{u \in \mathcal{A} : d(u, u(t)) = d_{\mathcal{A}}(u(t))\}$ erzwingen.
\item \emph{Approximative Prox-Regularität:} Herstellung einer modifizierten Reichweite nach Entfernung einer meageren Menge (Lebesgue-Null): Der Attraktor kann auf einer vernachlässigbaren Menge keine positive Reichweite haben, sie aber generisch beibehalten. Resultate von Lytchak~\cite{Lytchak2024} über die Struktur von Mengen positiver Reichweite legen nahe, dass eine solche Zerlegung durchführbar sein könnte.
\item \emph{Schwache Selektionsfunktionen:} Ersetzung starker Lipschitz-Projektionen durch Selektionsfunktionen beschränkter Hausdorff-Variation, die unter schwächeren Regularitätsannahmen existieren.
\item \emph{Geometrisches Euler-System:} Konstruktion einer kohärenten Hierarchie endlicher Galerkin-Approximationen $\mathcal{A}_N \to \mathcal{A}$ als ,,Reichweiten-Surrogat''. Für beliebige projizierte Attraktor-Punktwolken darf positive Reichweite nicht vorausgesetzt werden; der benötigte Input ist stattdessen ein verifizierter oder regularisierter Skeleton-Zweig mit kontrollierter Reichweite/Prox-Regularität und Hausdorff- oder Mosco-Konvergenz.
\end{enumerate}
\end{remark}

% ------------------------------------------------------------------
\subsection{Paper-übergreifende Verbindungen}
% ------------------------------------------------------------------

\begin{remark}[DFC aus Turbulenz als topologische Trichterheuristik für TLL]
\label{rem:dfc-funnel}
Die Abwärts-Flussbedingung (Downhill Flux Condition, DFC) aus dem Begleitpaper zur Turbulenz~\cite{FST-NS2026} formuliert, dass Energie ,,abwärts'' durch den Inertialbereich fließt. Übersetzt in die Attraktorgeometrie legt dies ein topologisches Trichterbild nahe: DFC kann die Annäherungsrichtung beschränken und dadurch eine log-Lipschitz-Rate unterstützen. Dies ist eine physikalische Heuristik für TLL, kein Ersatz für die Röhren-Projektionszweig-Abschätzung aus Theorem~\ref{thm:squeezing-tll}.
\end{remark}


% ===================================================================
% KI-Offenlegung
% ===================================================================
\subsection*{KI-Offenlegung}
Bei der Erstellung dieses Manuskripts wurden KI-gestützte Werkzeuge
(Claude, Anthropic) in unterstützender Funktion eingesetzt,
insbesondere für Literaturrecherche, Textstrukturierung und
Formulierungshilfe.  Die konzeptionelle Gestaltung, die
wissenschaftliche Argumentation und die Verantwortung für den
gesamten Inhalt liegen ausschließlich beim Autor.


\begin{thebibliography}{99}

\bibitem{AmbrosioGigliSavare2008}
L.~Ambrosio, N.~Gigli, and G.~Savar\'e,
\emph{Gradient Flows in Metric Spaces and in the Space of
Probability Measures},
2nd ed., Lectures in Mathematics ETH Z\"urich,
Birkh\"auser, Basel, 2008.

\bibitem{AmbrosioFuscoPallara2000}
L.~Ambrosio, N.~Fusco, and D.~Pallara,
\emph{Functions of Bounded Variation and Free Discontinuity Problems},
Oxford Mathematical Monographs, Clarendon Press, Oxford, 2000.

\bibitem{BabinVishik1992}
A.~V.~Babin and M.~I.~Vishik,
\emph{Attractors of Evolution Equations},
Studies in Mathematics and its Applications, vol.~25,
North-Holland, Amsterdam, 1992.

\bibitem{BaeCannone2016}
H.~Bae and M.~Cannone,
\emph{Log-Lipschitz regularity of the 3D Navier--Stokes equations},
Nonlinear Anal.\ \textbf{137} (2016), 222--245.

\bibitem{Chistyakov2004}
V.~V.~Chistyakov,
\emph{Selections of bounded variation},
J.~Appl.\ Anal.\ \textbf{10} (2004), no.~1, 1--82.

\bibitem{DeBlasiMyjak1993}
F.~S.~De~Blasi and J.~Myjak,
\emph{Ambiguous loci of the nearest point mapping in Banach spaces},
Arch.\ Math.\ \textbf{61} (1993), no.~4, 377--384.

\bibitem{DontchevRockafellar2009}
A.~L.~Dontchev and R.~T.~Rockafellar,
\emph{Implicit Functions and Solution Mappings: A View from
Variational Analysis},
Springer Monographs in Mathematics, Springer, New York, 2009.

\bibitem{ConstantinFoiasTemam1985}
P.~Constantin, C.~Foias, and R.~Temam,
\emph{Attractors representing turbulent flows},
Mem.\ Amer.\ Math.\ Soc.\ \textbf{53} (1985), no.~314.

\bibitem{EdenFoiasNicolaenkoTemam1994}
A.~Eden, C.~Foias, B.~Nicolaenko, and R.~Temam,
\emph{Exponential Attractors for Dissipative Evolution Equations},
Research in Applied Mathematics, vol.~37,
Wiley--Masson, Paris, 1994.

\bibitem{Federer1959}
H.~Federer,
\emph{Curvature measures},
Trans.\ Amer.\ Math.\ Soc.\ \textbf{93} (1959), 418--491.

\bibitem{FoiasManleyRosaTemam2001}
C.~Foias, O.~Manley, R.~Rosa, and R.~Temam,
\emph{Navier--Stokes Equations and Turbulence},
Encyclopedia of Mathematics and its Applications, vol.~83,
Cambridge University Press, 2001.

\bibitem{FST-NS2026}
L.~Geiger,
\emph{A Thermodynamic Obstruction to Finite-Time Blow-Up in the 3D
Navier--Stokes Equations (Conditional Framework)},
Zenodo, 2026.  DOI:\@ 10.5281/zenodo.19087449.

\bibitem{IlyinKalantarovZelik2025}
A.~Ilyin, V.~Kalantarov, and S.~Zelik,
\emph{Attractors and their dimensions for the 3D Fractional
Navier--Stokes--Voigt Equations},
arXiv preprint arXiv:2511.04783, 2025.

\bibitem{Lytchak2024}
A.~Lytchak,
\emph{A note on subsets of positive reach},
Math.\ Nachr.\ \textbf{297} (2024), no.~3, 932--942.

\bibitem{MalletParetSell1988}
J.~Mallet-Paret and G.~R. Sell,
\emph{Inertial manifolds for reaction diffusion equations in higher
space dimensions},
J.~Amer.\ Math.\ Soc.\ \textbf{1} (1988), no.~4, 805--866.

\bibitem{Mane1981}
R.~Man\'e,
\emph{On the dimension of the compact invariant sets of certain
non-linear maps},
Lecture Notes in Mathematics, vol.~898, Springer, Berlin, 1981,
pp.~230--242.

\bibitem{RockafellarWets1998}
R.~T.~Rockafellar and R.~J-B Wets,
\emph{Variational Analysis},
Grundlehren der mathematischen Wissenschaften, vol.~317,
Springer, Berlin, 1998.

\bibitem{PoliquinRockafellarThibault2000}
R.~A.~Poliquin, R.~T.~Rockafellar, and L.~Thibault,
\emph{Local differentiability of distance functions},
Trans.\ Amer.\ Math.\ Soc.\ \textbf{352} (2000), no.~11, 5231--5249.

\bibitem{Saks1937}
S.~Saks,
\emph{Theory of the Integral},
2nd revised ed., Monografie Matematyczne, vol.~7,
G.~E.~Stechert \& Co., New York, 1937.

\bibitem{Temam1997}
R.~Temam,
\emph{Infinite-Dimensional Dynamical Systems in Mechanics and Physics},
2nd ed., Applied Mathematical Sciences, vol.~68,
Springer, New York, 1997.

\bibitem{Zamfirescu1990}
T.~Zamfirescu,
\emph{The nearest point mapping is single valued nearly everywhere},
Arch.\ Math.\ \textbf{54} (1990), no.~6, 563--566.

\bibitem{FoiasTemam1979}
C.~Foias and R.~Temam,
\emph{Some analytic and geometric properties of the solutions of the
evolution Navier--Stokes equations},
J.\ Math.\ Pures Appl.\ \textbf{58} (1979), no.~3, 339--368.

\bibitem{Romanov2000}
A.~V.~Romanov,
\emph{Sharp estimates of the dimension of inertial manifolds for
nonlinear parabolic equations},
Izv.\ Ross.\ Akad.\ Nauk Ser.\ Mat.\ \textbf{57} (1993), 36--52;
English transl.\ in Russian Acad.\ Sci.\ Izv.\ Math.

\end{thebibliography}

\end{document}



